% concatenated from ArP_SK2_v6_arvix.tex
\documentclass[12pt]{article}

% Language setting
% Replace `english' with e.g. `spanish' to change the document language
\usepackage[english]{babel}

% Set page size and margins
% Replace `letterpaper' with`a4paper' for UK/EU standard size
\usepackage[a4paper,top=2cm,bottom=2cm,left=3cm,right=3cm,marginparwidth=1.75cm]{geometry}

% Useful packages
\usepackage{amsmath}
\usepackage{graphicx}
\usepackage{yfonts}
\usepackage{amssymb}
\usepackage{amsthm,bm}
\usepackage{comment}
\usepackage{bbm}
\usepackage{mathtools}
\newtheorem{thm}{Theorem}[section]
\newtheorem{co}[thm]{Corollary}
\newtheorem{lem}[thm]{Lemma}
\newtheorem{conj}[thm]{Conjecture}
\newtheorem{assumption}[thm]{Assumption}
\newenvironment{assum}{\begin{assumption}\rm}{\end{assumption}}
\newtheorem{pr}[thm]{Proposition}
\newtheorem{crit}[thm]{Criterion}
\newtheorem{prob}[thm]{Problem}
\newtheorem{definition}[thm]{Definition}
\newenvironment{de}{\begin{definition}\rm}{\end{definition}}
\newtheorem{example}[thm]{Example}
\newenvironment{exmp}{\begin{example}\rm}{\end{example}}
\newtheorem{remark}[thm]{Remark}
\newenvironment{rem}{\begin{remark}\rm}{\end{remark}}
\newenvironment{claim}[1]{\par\noindent\underline{Claim:}\space#1}{}
\newenvironment{claimproof}[1]{\par\noindent\underline{Proof:}\space#1}{\hfill $\blacksquare$}
\usepackage[colorlinks=true, allcolors=blue]{hyperref}
\renewcommand\thesection{\arabic{section}}
\renewcommand\thesubsection{\thesection.\arabic{subsection}}

\newcommand{\E}{\mathbb{E}}
\newcommand{\EX}{{\mathbb{E}}}
\newcommand{\PX}{{\mathbb{P}}}
\newcommand{\R}{\mathbb R}
\newcommand{\eps}{\varepsilon}
\newenvironment{keywords}
  {\par\noindent\textbf{Keywords: }}
  {\par}

\DeclareMathOperator{\real}{\mathbb{R}}
\DeclareMathOperator{\borel}{\mathcal{B}}

\title{Second-Order Schalkwijk-Kailath Coding for Autoregressive Gaussian Channels}
\author{\footnotesize \centering\begin{tabular}{ccc}
Jun Su & Guangyue Han & Shlomo Shamai (Shitz)\\
The University of Hong Kong &The University of Hong Kong& Technion-Israel Institute of Technology\\
 email:  junsu@hku.hk& email:  ghan@hku.hk& email: sshlomo@ee.technion.ac.il\\
\end{tabular}}
\begin{document}
\maketitle

\begin{abstract}
We study communication with noiseless feedback over Gaussian channels with stationary autoregressive noise of arbitrary finite order. We introduce a class of Gaussian feedback coding schemes, called SK(2), in which the message process follows a second-order deterministic recursion, and derive a closed-form characterization of its maximal achievable rate under an average-power constraint. As a first-order benchmark, we formulate the SK(1) scheme and show that it provides the branch-complete and sign-consistent reformulation of Butman's equal-energy linear-feedback construction. The SK(2) coding scheme achieves feedback capacity for the additive white Gaussian noise channel and stationary AR(1) Gaussian channels. For certain stationary AR(2) Gaussian channels, genuinely second-order SK(2) scheme strictly outperforms SK(1); for the subclass obtained by interleaving two independent AR(1) noise processes, SK(2) also achieves feedback capacity. These results show that first-order SK/Butman coding scheme is not universally optimal beyond first-order autoregressive noise and disprove the corrected form of Butman's conjecture.



%\textcolor{red}{We study feedback communication over stationary AR($p$) Gaussian channels and propose a Gaussian random coding scheme that generalizes the celebrated Schalkwijk--Kailath (SK) coding scheme. The proposed coding scheme, termed the SK(2) coding scheme, is based on a second-order deterministic message recursion and admits a closed-form characterization of its maximal achievable rate. We also formulate the SK(1) coding scheme as the complete sign-consistent version of Butman's equal-energy first-order linear-feedback scheme. For certain AR(2) Gaussian channels, we show that the SK(2) coding scheme strictly outperforms the SK(1), or corrected Butman, coding scheme. Consequently, the first-order SK/Butman scheme is not universally optimal for AR($p$) Gaussian channels with $p>1$, disproving the corrected form of Butman's conjecture.}

%We propose a random coding scheme for AR($p$) Gaussian channels that generalizes the celebrated Schalkwijk-Kailath (SK) coding scheme. This constructive coding scheme, termed the SK(2) coding scheme, yields a closed-form characterization for the corresponding achievable rate. Among many others, this result shows that the SK coding scheme is not universally optimal, thereby disproving the conjecture proposed by Butman in \cite{butman1976linear}.
\end{abstract}

\begin{keywords}
autoregressive Gaussian channels, feedback capacity, linear-feedback coding, Schalkwijk--Kailath coding, colored Gaussian noise
\end{keywords}

\section{Introduction} \label{sec:introduction}
We consider the following additive colored Gaussian noise (ACGN) channel with noiseless feedback:
\begin{equation}\label{Channel-ACGN}
   Y_i = X_i(M,Y_1^{i-1}) + Z_i,\quad i=1,2,\dots,
\end{equation}
where $M$ denotes the transmitted message, and $\{Z_i\}_{i=1}^\infty$ is a zero-mean stationary Gaussian noise process independent of $M$. The channel input $X_i$ at time $i$ may causally depend on the message $M$ and the previous channel outputs $Y_1^{i-1}\triangleq\{Y_1,\ldots,Y_{i-1}  \}$. The channel input sequence $\{X_i\}_{i=1}^\infty$ is subject to an average power constraint $P>0$, namely,
$$
   \frac{1}{n}\sum_{i=1}^n\E[X^2_i(M,Y_1^{i-1})] \le P ,\qquad n\ge1.
$$
Let $C_{\mathrm{FB}}(P)$ denote the capacity of the ACGN channel \eqref{Channel-ACGN}, commonly referred to as the {\em Gaussian feedback capacity}.

It is well-known that, without feedback, the capacity of the ACGN channel \eqref{Channel-ACGN}, denoted by $C(P)$, can be determined by the water-filling method \cite{Shannon1949}. In the special case where the noise process is white, i.e., $Z_i=W_i$ with $\{W_i\}$ a white Gaussian process with zero mean and unit variance, the channel reduces to the additive white Gaussian noise (AWGN) channel, whose capacity was first determined by Shannon \cite{shannon1948mathematical,Shannon1949}; see also \cite{wolfowitz1964,wolfowitz1968note}. More precisely, Shannon \cite{shannon1956zero} proved that for this channel, feedback does not increase capacity; that is, the feedback capacity $C_{\mathrm{FB}}(P)$ coincides with the nonfeedback capacity $C(P)$; see also \cite{kadota1971capacity,kadota1971mutual}. Nevertheless, feedback can greatly simplify coding and improve reliability. A significant later result, due to Schalkwijk and Kailath \cite{schalkwijk1966coding2, schalkwijk1966coding}, showed that a simple linear feedback coding scheme, commonly referred to as the SK coding scheme, achieves the capacity of the AWGN channel.

For the ACGN channel \eqref{Channel-ACGN} when the noise process is not white, feedback may increase capacity, and the optimal coding structure is substantially more delicate. Butman \cite{butman1969} formulated linear feedback communication systems for general ACGN channels and derived the corresponding optimal linear signaling structure. For stationary first-order autoregressive, or AR(1), Gaussian channels under average power constraint \(P\), Butman's equal-energy scheme gives a closed-form achievable rate. Wolfowitz \cite{wolfowitz1975signalling} provided an alternative derivation for such rate and proved optimality within a class of linear signaling schemes. Tiernan and Schalkwijk \cite{Tiernan1974} established upper bounds on $C_\mathrm{FB}(P)$ for AR(1) Gaussian channels with noiseless feedback, while Ozarow \cite{ozarow1990upper,ozarow1990random} studied upper bounds on $C_\mathrm{FB}(P)$ and random coding for ACGN channels with feedback. In particular, Ozarow \cite{ozarow1990random} showed that feedback strictly increases capacity for nonwhite stationary Gaussian noise, while traditional random coding is not generally capacity-achieving in the feedback setting.

Butman subsequently extended the equal-energy construction to general $p$-th-order autoregressive, or AR($p$), Gaussian channels and conjectured that the resulting lower bound is tight \cite{butman1976linear}. As we explain in Section~\ref{subsec:relation_to_butman}, however, the AR($p$) formula in \cite{butman1976linear} contains a sign issue: the transition from the absolute-value recursion to the final polynomial expression requires additional sign-compatibility and branch conditions. As detailed later on, we will resolve this issue and present the corrected form of Butman's scheme.

Several later works studied feedback capacity for the ACGN channel \eqref{Channel-ACGN} from dynamic-programming, state-space, and variational viewpoints. Yang et al. \cite{yang2005feedback,yang2007feedback} developed dynamic-programming formulations for finite-state machine channels and power-constrained Gaussian channels with memory. For the first-order moving-average, or MA(1), Gaussian channel, Kim \cite{Kim2006} showed that a modified SK coding scheme achieves the feedback capacity and obtained a closed-form expression for $C_\mathrm{FB}(P)$. Subsequently, Kim \cite{kim2009feedback} gave a variational characterization of the feedback capacity of stationary Gaussian channels. Related refinements and extensions can be found in \cite{liu2018feedback,sabag2023,su2024feedback}. In particular, the first-order autoregressive moving-average, or ARMA(1,1), Gaussian channel is now well understood. Kim \cite[Theorem 6.1]{kim2009feedback} also stated that a $k$-dimensional variant of the SK coding scheme achieves feedback capacity for any ARMA noise of order $k$. However, Derpich and \O stergaard \cite{derpich2022comments} identified a gap in the proof of a key intermediate result \cite[Corollary 4.4]{kim2009feedback}. Consequently, the proof of Theorem 6.1 in \cite{kim2009feedback} is incomplete, leaving the optimality of the proposed $k$-dimensional coding scheme open in general.

%Linear feedback schemes have also been studied for AWGN channels with noisy output feedback. Chance and Love \cite{chance2011concatenated} proposed a linear scheme that significantly outperforms the SK coding scheme for AWGN channels with AWGN noisy feedback, and Mishra et al. \cite{mishra2023linear} later used dynamic programming to optimize a class of sequential linear schemes. These works address noisy-feedback models and are therefore complementary to the noiseless-feedback AR($p$) Gaussian channels considered here.

%For the general $p$th-order autoregressive (AR($p$)) Gaussian channel under average power constraint $P$, where the noise process $\{Z_i\}$ in \eqref{Channel-ACGN} is a stationary AR$(p)$ Gaussian process, Butman \cite{butman1969} derived the optimal linear signaling structure with feedback and for stationary AR(1) Gaussian channels gave a closed-form lower bound on $C_\mathrm{FB}(P)$ (which coincides with $C_\mathrm{SK1}(P)$ in Theorem \ref{ISK1-thm} when $p=1$), achievable under equal signal energy allocation. Wolfowitz \cite{wolfowitz1975signalling} provided an alternative derivation of the same bound and proved its optimality within a class of linear signaling schemes.  This lower bound is in fact tight, as shown in \cite{kim2009feedback,yang2007feedback,su2024feedback}. Upper bounds on $C_{\mathrm{FB}}(P)$ of the AR($p$) Gaussian channel were further provided in \cite{Tiernan1974}.

 % More recently, Kim in \cite[Section 6]{kim2009feedback} showed that the SK coding scheme achieves the feedback capacity for the first-order autoregressive moving-average (ARMA(1,1)) Gaussian channel. Using a state-space approach, Kim further stated in \cite[Theorem 6.1]{kim2009feedback} that a $k$-dimensional variant of the SK coding scheme achieves the feedback capacity for any ARMA noise spectrum of order $k$. However, as noted in a recent paper \cite{derpich2022comments}, there is a gap in the proof of a key intermediate result \cite[Corollary 4.4]{kim2009feedback}. Consequently, the proof of Theorem 6.1 in \cite{kim2009feedback} is incomplete, leaving the optimality of the scheme as a conjecture.


%\footnote{As noted in a recent paper \cite{derpich2022comments}, there is a gap in the proof of a key intermediate result \cite[Corollary 4.4]{kim2009feedback}; consequently, the proof of Theorem 6.1 in \cite{kim2009feedback} is incomplete, leaving the optimality of the scheme as a conjecture.}

In this paper, we focus on the stationary AR($p$) Gaussian channel
\begin{equation}\label{channel-arp}
Y_i = X_i + Z_i, \quad i\ge 1,
\end{equation}
where the noise process $\{Z_i\}$ is a real-valued AR($p$) Gaussian process. Its power spectral density \cite{liu2018feedback} is
\begin{equation}\label{sdf-arp}
S_Z(e^{i\theta}) = \vert H_Z(e^{i\theta}) \vert^2 = \frac{1}{\vert L_Z(e^{i\theta}) \vert^2} = \frac{1}{\vert\prod_{k=1}^{p}(1+\beta_ke^{i\theta})\vert^2},
\end{equation}
where $|\beta_k|<1$, $k=1,\ldots,p$, and the parameters are chosen so that \(L_Z\) has real coefficients. As discussed above, the literature on optimal coding for stationary ARMA Gaussian channels is extensive. However, to the best of our knowledge, constructing the optimal coding scheme and  analytically determining the corresponding feedback capacity, remain an open problem for all ARMA orders beyond the ARMA(1,1) case. 

We first formulate the SK(1) coding scheme, which gives a branch-complete and sign-consistent version of Butman's equal-energy first-order linear-feedback scheme \cite{butman1976linear}. With this correction, Butman's conjecture can be stated as the optimality of the SK(1) coding scheme:
$$
    C_{\mathrm{FB}}(P)=C_{\mathrm{SK1}}(P),
$$
where $C_{\mathrm{SK1}}(P)$ is the maximal achievable rate of the SK(1) coding scheme. The main contribution of this work is a closed-form expression for the maximum achievable rate of a Gaussian random coding scheme, termed the SK(2) coding scheme, over the stationary AR($p$) channel (\ref{channel-arp}). This rate, denoted by $C_\mathrm{SK2}(P)$, provides a lower bound on $C_\mathrm{FB}(P)$. For certain AR(2) Gaussian channels, we show that the SK(2) coding scheme can strictly outperform the SK(1) coding scheme and hence
$$
    C_{\mathrm{FB}}(P)\ge C_{\mathrm{SK2}}(P)>C_{\mathrm{SK1}}(P),
$$
which disproves the corrected form of Butman's conjecture.


The remainder of the paper is organized as follows. In Section~\ref{sec:the_classical_schalkwijk_kailath_coding}, we review the classical SK coding scheme. In Section~\ref{subsection:the_sk_1_coding_scheme}, we introduce the SK(1) coding scheme and derive its maximal achievable rate in Theorem~\ref{ISK1-thm}. We then clarify its relation to Butman's equal-energy scheme and reformulate Butman's conjecture in terms of \(C_{\mathrm{SK1}}(P)\) in Section~\ref{subsec:relation_to_butman}. Section~\ref{sec:sk2-coding-scheme} contains the main contributions:
\begin{itemize}
\item[$\bullet$] We propose the SK(2) coding scheme for AR($p$) Gaussian channels, a second-order generalization of the SK(1) coding scheme, and derive its maximal achievable rate (Theorem~\ref{main-theorem-sk2}, Corollary~\ref{main-col-sk2:label}).
\item[$\bullet$]  For AWGN and AR(1) Gaussian channels, we show that the SK(2) coding scheme achieves feedback capacity (Corollary~\ref{co-awgn-sk2:label} and Corollary~\ref{co-ar1-sk2:label}).
\item[$\bullet$] For stationary AR(2) Gaussian channels, numerical examples demonstrate that the SK(2) coding scheme can strictly outperform the SK(1) coding scheme. We further identify an interleaved AR(1) family for which the SK(2) coding scheme achieves feedback capacity and strictly outperforms the SK(1) coding scheme, thereby disproving the corrected form of Butman's conjecture (Proposition~\ref{prop:interleaved-AR1}).
%For AR(2) Gaussian channels, we show that the SK(2) coding scheme can strictly outperform the SK(1) coding scheme, or corrected Butman's coding scheme. This disproves the corrected form of Butman's conjecture \cite{butman1976linear} (Example~\ref{example-ar2:label}).
\end{itemize}

%\section{Notation and Terminologies}\label{sec:notation_and_terminologies}
%We use $(\Omega,\mathcal{F},\mathbb{P})$ to denote the underlying probability space, and $\E$ to denote the expectation with respect to the probability measure $\mathbb{P}$. Let $X,Y$ be random variables defined on the probability space $(\Omega,\mathcal{F},\mathbb{P})$, which will be used to illustrate most of the notions and facts in this section (the same notations may have different connotations in other sections). In this work, a random variable may be real-valued (with a probability density function) or vector-valued (more precisely, taking values in $\mathbb{R}^\infty$ or $\mathbb{R}^n$). We use $\mu_{X}$ and $\mu_Y$ to denote the induced probability distributions of $X$ and $Y$, respectively. Their product distribution is denoted by $\mu_X\times\mu_Y$, and their joint distribution is $\mu_{XY}$. Throughout the remainder of this paper, unless otherwise specified, $\{W_i\}$ denotes a white Gaussian process with zero mean and unit variance.

% The {\it mutual information} $I(X;Y)$ between two random variables $X,Y$ is defined as
% { \begin{equation*} \label{definition-mutual-information}
% I(X; Y)=\begin{cases}
% \EX\left[\log \frac{d \mu_{XY}}{d \mu_{X} \times \mu_{Y}}(X, Y)\right], & \mbox{ if } \frac{d \mu_{X {Y}}}{d \mu_{X} \times \mu_{Y}} \mbox{ exists},\\
% \infty, & \mbox{otherwise},
% \end{cases}
% \end{equation*}}
% \!\!\!where $d \mu_{X{Y}}/d \mu_{X} \times \mu_{Y}$ denotes the Radon-Nikodym derivative of $\mu_{XY}$ with respect to $\mu_{X} \times \mu_{Y}$.

%We use $\delta_{ij}$ to denote the Kronecker delta function. 
%For $a\in\mathbb{C}$, $|a|$ denotes its modulus; for $a\in\mathbb{C}^n$, $\|a\|$ denotes its $\ell^2$-norm. For an $m\times n$ matrix $A$, we use $A^\ast$ to denote its conjugate transpose (or $A^T$ if $A$ is real). If $A$ is square, $|A|$ denotes its determinant, with the convention that $|A| = 1$ for an empty matrix (i.e., $m=n=0$).


\section{The SK(1) Coding Scheme}
It is well known that the SK coding scheme achieves feedback capacity for the AWGN channel under an average power constraint $P$. In this section, we extract its core structure into a first-order deterministic message recursion, termed the SK(1) coding scheme. We first recall the classical SK coding scheme as motivation. We then define the SK(1) coding scheme for stationary AR($p$) Gaussian channels \eqref{channel-arp} and derive its maximal achievable rate. Finally, we relate the SK(1) coding scheme to Butman's equal-energy linear-feedback coding scheme \cite{butman1976linear}, point out some technical issues in Butman's treatment, and reformulate Butman's conjecture in terms of the corrected Butman's lower bound.


\subsection{The Schalkwijk-Kailath Scheme}\label{sec:the_classical_schalkwijk_kailath_coding}
In the spirit of the exposition in \cite[Section 17.1.1]{el2011network}, the classical SK coding scheme can be reformulated as follows. Consider the discrete-time AWGN channel
\begin{equation}\label{AWGN-channel}
Y_i=X_i+W_i,\quad i \ge 0,
\end{equation}
where $\{W_i\}$ is a white Gaussian process with zero mean and unit variance and $\{X_i\}$ is subject to the average power constraint $P$, i.e., 
$$
\frac{1}{n}\sum_{i=0}^{n-1}\E[X_i^2]\le  P.
$$
The encoding scheme proceeds as follows. At time $i=0$,  the encoder transmits a symbol $X_0$ chosen arbitrarily from the interval $[-\sqrt{P},\sqrt{P}]$,  yielding the output $Y_0 = X_0+W_0$. Because of noiseless feedback, the encoder learns $Y_0$ and can compute the noise realization $W_0 = Y_0 - X_0$. At time $i=1$, the encoder transmits $X_1 = \gamma_1W_0$, where $\gamma_1=\sqrt{P}$ ensures that $\E[|X_1|^2] =P$. For each subsequent time $i\ge 2$, the encoder computes the minimum mean-square error (MMSE) estimate $\E[W_0\mid Y_1^{i-1}]$ based on all previous channel outputs $Y_1^{i-1}$ available via feedback. It then transmits 
\begin{equation}\label{SK-coding-rule}
X_i = \gamma_i (W_0 - \E[W_0\mid Y_1^{i-1}]),
\end{equation}
where $\gamma_i$ is chosen such that $\E[|X_i|^2] = P$ for all $i$. The SK coding scheme, given in (\ref{SK-coding-rule}), admits a recursive interpretation in which the transmitter repeatedly refines the receiver's knowledge of the initial noise $W_0$, or equivalently, the original message $X_0$, across successive uses of the channel. Specifically, the encoding rule can be written as
\begin{equation*}
\begin{aligned}
X_i &=  \gamma_i (W_0 - \E[W_0\mid Y_1^{i-1}])\\
&= \frac{\gamma_i}{\gamma_{i-1}} \left(X_{i-1} -\E[X_{i-1}\mid Y_1^{i-1}]\right)
\end{aligned}
\end{equation*}
so that each transmission corrects the receiver's estimation error incurred in the previous transmission. 

Furthermore, define the process $V_i \triangleq \gamma_i W_0$. Since $W_0$ is independent of the forward noise $\{W_i\}_{i\ge 1}$, $\{V_i\}$ may be regarded as the underlying {\it message process}. Under this representation, the SK rule takes the following form
\begin{equation}\label{SK-coding-rule-X}
X_i = V_i - \E[V_i\mid Y_1^{i-1}],
\end{equation}
where $\{V_i\}$ follows the deterministic recursion:
\begin{equation}\label{SK-coding-rule-V}
V_i = \frac{\gamma_i}{\gamma_{i-1}} V_{i-1},\quad i\ge 2
\end{equation}
with initial condition $V_1=\sqrt{P}W_0$. It has been shown \cite{schalkwijk1966coding2,schalkwijk1966coding,el2011network,durrett2019probability} that for any such sequence $\{\gamma_i\}_{i\ge1}$, this SK coding scheme (\ref{SK-coding-rule}), or equivalently, the recursive form (\ref{SK-coding-rule-X}) and (\ref{SK-coding-rule-V}), achieves the feedback capacity of the AWGN channel.


The SK coding scheme described above is designed for the AWGN channel \eqref{AWGN-channel}, where the noise samples are uncorrelated. For the general ACGN channel \eqref{Channel-ACGN}, the above recursive representation motivates a natural first-order generalization. We now formalize this generalization as the SK(1) coding scheme. Its relation to Butman's equal-energy coding scheme will be discussed after the corresponding maximal achievable rate characterization in Theorem~\ref{ISK1-thm}.

\subsection{The SK(1) Coding Scheme}\label{subsection:the_sk_1_coding_scheme}
We now introduce a first-order specialization of the SK coding scheme, which we term the SK(1) coding scheme. The name reflects that the underlying message process follows a first-order deterministic recursion. It is obtained from the recursive form in (\ref{SK-coding-rule-X}) and (\ref{SK-coding-rule-V}) by imposing the simple ratio condition
\begin{equation}\label{sk1-ratio}
\frac{\gamma_i}{\gamma_{i-1}}=\gamma\quad\mbox{for all}~i,
\end{equation}
where $|\gamma|>1$. Under this condition, the message process $V=\{V_i\}$ follows a first-order deterministic recursion. Formally, the SK(1) coding scheme is defined by
\begin{equation}\label{SK1-coding-X}
X_i = V_i - \E[V_i\mid Y_1^{i-1}]
\end{equation}
with the message process $\{V_i\}$ such that 
\begin{equation}\label{SK1-coding-message}
\begin{cases}
V_{i+1} = \gamma V_i, & i \ge 1,\\
V_1 = U,
\end{cases}
\end{equation}
where $U$ is a standard Gaussian random variable, independent of the noise process $Z=\{Z_i\}$. We define the {\em SK(1) information rate capacity}, or simply, {\em SK(1) capacity} under power constraint \(P\) as
\begin{equation}\label{feedback-capacity-SK1}
C_{\mathrm{SK1}}(P) = \sup_{(V,X)} \bar{I}(V;Y),
\end{equation}
where $\bar{I}(V;Y)$ denotes the mutual information rate between the processes $\{V_i\}$ and \(\{Y_i\}\). The supremum is taken over all pairs \((V,X)\) satisfying \eqref{SK1-coding-X}, \eqref{SK1-coding-message} and the asymptotic average power
constraint
\begin{equation}\label{power-constraint-sk1}
   \varlimsup_{n\to\infty}\frac{1}{n}\sum_{i=1}^n\E[X_i^2]\le P.
\end{equation}
Thus, \(C_{\mathrm{SK1}}(P)\) is the maximal achievable information rate within the SK(1) coding class. In particular,
$$
    C_{\mathrm{SK1}}(P)\le C_{\mathrm{FB}}(P).
$$
\begin{remark}
Note that the SK(1) coding scheme, which can be directly transformed into the original SK coding, has been established as capacity-achieving for stationary ARMA(1,1) Gaussian channels with feedback \cite{kim2009feedback}.
\end{remark}

We begin by stating a key lemma, the proof of which is deferred to Appendix \ref{sec:proof_of_lemma_ref_lemma1}. This lemma characterizes the quadratic form associated with the inverse of a structured matrix.
\begin{lem}\label{lemma1}
Let $\{u_k;k=1,2,...,m\}\subset \mathbb{R}^{n}$ for $n,m \ge 1$ and let $S$ be an $n\times n$ real symmetric positive-definite matrix. Define $U = \begin{bmatrix} u_1 & u_2 & \cdots & u_m  \end{bmatrix}$ and $A = I_m +U^T S^{-1} U$, where $I_m$ is the $m\times m$ identity matrix. Then for any $1\le i,j\le m$,
$$
u_i^T (S+UU^T)^{-1}u_j= \delta_{ij} + (-1)^{i+j+1} \frac{ M_{ji}}{\det(A)}
$$
where $\delta_{ij}$ is the Kronecker delta function, $M_{ij}$ is the $(i, j)$-minor of $A$ and we have adopted the standard notation $\det(A)$ for the determinant of $A$. In particular, when $m=1$,
$$
u_1^T (S+u_1u_1^T)^{-1}u_1=  1-\frac{1}{A}.
$$
\end{lem}

The following theorem gives a closed-form expression for the SK(1) capacity $C_{\mathrm{SK1}}(P)$ over stationary AR$(p)$ Gaussian channels, which provides a computable lower bound on the feedback capacity $C_\mathrm{FB}(P)$.
\begin{thm}\label{ISK1-thm}
Suppose the noise $\{ Z_i \}$ is an $AR(p)$ Gaussian process with parameters $\beta_k$, $\vert \beta_k\vert<1$ for all $k=1,2,...,p$, namely, it has the power spectral density
\begin{equation*}
S_Z(e^{i\theta}) = \vert H_Z(e^{i\theta}) \vert^2 = \frac{1}{\vert L_Z(e^{i\theta}) \vert^2} = \frac{1}{\vert\prod_{k=1}^{p}(1+\beta_ke^{i\theta})\vert^2}.
\end{equation*}
Then, we have
\begin{equation}\label{I-SK1-thm}
C_\mathrm{SK1}(P)= \max_{\gamma} \log  |\gamma|,
\end{equation}
where the maximum is taken over all $\gamma\in\mathbb{R}$ with $|\gamma|>1$ satisfying the power constraint 
\begin{equation}\label{power-eqn-sk1-thm}
   (\gamma^2-1)H^2_Z(\gamma^{-1}) = P.
\end{equation}
\end{thm}
\begin{remark}\label{remark1:label}
For the special cases $p=0$ and $p=1$, corresponding to the AWGN channel and AR(1) Gaussian channel, respectively, the feedback capacity $C_{\mathrm{FB}}(P)$ under power constraint $P$ is already known from the existing literature (see \cite{kim2009feedback,yang2007feedback} and \cite[Theorem 5.2]{su2024feedback}). It follows directly from these results that $C_\mathrm{SK1}(P)=C_{\mathrm{FB}}(P)$ for both cases. Therefore, the SK(1) coding scheme is optimal for the AWGN and AR(1) Gaussian channels.
\end{remark}

\begin{proof}
The proof follows a similar argument to that of our main result Theorem~\ref{main-theorem-sk2} and is therefore deferred to Appendix \ref{sec:proof_of_theorem_ref_isk1_thm}.
\end{proof}

%Theorem~\ref{ISK1-thm} gives a rigorous first-order benchmark for stationary AR($p$) Gaussian channels \eqref{channel-arp}. We next show that this benchmark is also the natural corrected form of Butman's equal-energy lower bound \cite{butman1976linear}.

%To conclude this section, we have introduced the $\mathrm{SK}(1)$ coding scheme and derived an explicit expression for the maximal rate it achieves, as presented in Theorem~\ref{ISK1-thm}. This result extends the formula for the SK(1) coding scheme from AR(1) to AR($p$) Gaussian channels. In the following section, we introduce the SK(2) coding scheme, wherein the underlying message process follows a second-order deterministic recursion, thereby generalizing the $\mathrm{SK}(1)$ coding scheme.

\subsection{Butman's Equal-Energy Coding Scheme}\label{subsec:relation_to_butman}
In this subsection, we review Butman's equal-energy linear-feedback coding scheme in the notation of this paper. We first recall the AR(1) case, where Butman's construction gives the known feedback capacity. We then explain two issues in the AR($p$) extension: a sign-compatibility issue between the recursive sign selection and the equal-ratio specialization, and a branch issue in removing the absolute value from Butman's recursion. The detailed AR(2) counterexamples are deferred to Appendix~\ref{sec:example_of_ar_2_gaussian_channel}. Finally, we formulate the corrected branch-complete Butman rate and show that it coincides with the SK(1) capacity.

Butman extended the SK coding scheme from the AWGN channel to the ACGN channel \eqref{Channel-ACGN} (see \cite{butman1969,butman1976linear}). More precisely, the linear-feedback signaling rule can be formulated as
$$
    X_i=g_i\bigl(U-\mathbb E[U\mid Y_1^{i-1}]\bigr),
$$
where $U$ is a Gaussian message variable independent of the noise, and \(\{g_i\}\subset\mathbb R\) is chosen to satisfy the average power constraint $P$. For the stationary AR(1) Gaussian channel with coefficient \(\beta_1\), Butman imposed an equal-energy allocation and a recursive sign selection. Defining
\begin{equation}
    x_i=\left|\frac{g_i}{g_{i-1}}\right|,
\end{equation}
and choosing the sign of \(g_i\) such that
\begin{equation}
    \operatorname{sign}(g_i)=\operatorname{sign}(\beta_1 g_{i-1}),
\end{equation}
Butman obtained
\[
    x_{i+1}^2 = 1+P\left(1+\frac{|\beta_1|}{x_i}\right)^2 .
\]
Under the equal-energy allocation \(x_i=x>1\), this gives
\begin{equation}\label{butman-original-rate:label}
    R^{(1)}_{ee}(P)=\log x,
\end{equation}
where $x$ is the unique positive root of the following polynomial equation
\[
    P=\frac{x^2-1}{(1+|\beta_1|x^{-1})^2}.
\]
This expression agrees with the known feedback capacity of the AR(1) Gaussian channel \cite{kim2009feedback,yang2007feedback,sabag2023}.

Butman further proposed an extension of this construction to the stationary AR($p$) Gaussian channel \eqref{channel-arp}. In the notation of this paper, we write
\[
    L_Z(z)=\prod_{k=1}^{p}(1+\beta_k z)
          =1+\sum_{k=1}^{p}b_k z^k
\]
for some $\{b_k;k=1,2,\cdots,p\}\subset \mathbb{R}$, and define
\[
    B(x)=\sum_{k=1}^{p}b_kx^{-k}.
\]
After the recursive sign selection and the equal-energy allocation, Butman's derivation leads to the following absolute-value equation
\begin{equation}\label{butman-abs-eqn}
    x^2= 1+ P\left(1+|B(x)|\right)^2 .
\end{equation}
Thus, if the equation \eqref{butman-abs-eqn} is valid, the corresponding achievable rate is
$$
    R^\mathrm{abs}_{ee}(P)=\frac12\log x^2=\log|x|,
$$
for any real root $x$ satisfying \eqref{butman-abs-eqn} and $|x|>1$.

Some issue occurs in the step from \eqref{butman-abs-eqn} to Butman's final polynomial expression (\cite[Equation (49)]{butman1976linear}), which in the present notation corresponds to
\begin{equation}\label{butman-final-eqn-1}
    x^2= 1+ P\left(1-B(x)\right)^2
\end{equation}
and the corresponding achievable rate is 
\begin{equation}\label{Butman-ARp-mutual-info}
    R_{ee}^{(p)}(P) = \max \log|x|, 
\end{equation}
where the maximum is taken over all real roots $x$ of the equation \eqref{butman-final-eqn-1} satisfying $|x|>1$. This step removes the absolute value in \eqref{butman-abs-eqn} with a fixed sign, which is valid if $B(x)$ is negative.

Indeed, Butman noted that ``A little thought will reveal that the above choice of sign for \(g_n\), may not be the best. This problem deserves more investigation. However, if all the \(\alpha\) have a negative sign, the above choice is the best.'' Note that $\alpha_k$ in \cite{butman1976linear} corresponds to $-b_k$ in this paper. Hence, Butman's sufficient condition that all \(\alpha_k\) are negative corresponds to
$$
    b_k>0,\qquad k=1,\ldots,p .
$$
On the positive branch \(x>0\), this implies $B(x)>0$, and immediately $|B(x)|=B(x)$. Consequently, on this positive branch, the sign-consistent equation following from \eqref{butman-abs-eqn} is not \eqref{butman-final-eqn-1}, but rather
\begin{equation}\label{butman-final-eqn-2}
    x^2=1+P\left(1+B(x)\right)^2 .
\end{equation}
Thus, on this positive branch, the corresponding achievable rate is
$$
    R^{(+)}_{\mathrm{Butman}}(P)=\max \log x,
$$
where the maximum is taken over all real roots $x$ of \eqref{butman-final-eqn-2} satisfying $x>1$. Therefore, the transition from \eqref{butman-abs-eqn} to \eqref{butman-final-eqn-1} is not valid in general; the correct sign depends on the value of $B(x)$ at the selected root $x$.
%an explicit branch condition is required, and on the branch suggested by Butman's own sufficient condition the sign is opposite to that in \eqref{butman-final-eqn-1}.


There is a further compatibility issue behind \eqref{butman-abs-eqn}. In Butman's derivation for the AR($p$) case, the ratios $x_i=g_i/g_{i-1}$ are signed quantities. Thus, the equal-energy allocation $x_i=x$, \(|x|>1\), restricts the sign selection of \(\{g_i\}\). Indeed, let $s_i=\operatorname{sign}(g_i) \in \{1,-1\}$. If the equal-energy allocation $x_i=x$ is imposed, then
\begin{equation}\label{Butman-ee-cond}
    s_i=\operatorname{sign}(x)s_{i-1}.
\end{equation}
Thus, the sign sequence $\{s_i\}$ is constant when $x>0$ and alternating when $x<0$. Furthermore, let $r=|x|$ and the squared term in Butman's recursion (see \cite[Equation (39)]{butman1976linear}) can be written as 
\begin{equation}\label{Butman-recursive-term}
    \left(1+\sum_{k=1}^{p}b_k\frac{g_{i-k}}{g_i}\right)^2 = \left(1+s_i\sum_{k=1}^{p}b_k s_{i-k}r^{-k}\right)^2.
\end{equation}
Define
\[
    S_i(r)=\sum_{k=1}^{p} b_k s_{i-k} r^{-k}, 
\]
For fixed $\{s_{i-1},\cdots,s_{i-p}\}$, the two possible choices $s_i=1$ and $s_i=-1$ give $(1+S_i(r))^2$ and $(1-S_i(r))^2$, respectively. Therefore, to maximize \eqref{Butman-recursive-term}, when \(S_i(r)\neq0\), Butman's recursive sign choice is
\begin{equation}\label{Butman-sign-cond}
    s_i=\operatorname{sign}(S_i(r)),
\end{equation}
and the squared term \eqref{Butman-recursive-term} becomes
\begin{equation}\label{Butman-one-step-factor}
    \left(1+|S_i(r)|\right)^2.
\end{equation}
%These two requirements are compatible for $p=1$, but not necessarily for $p>1$.

On the other hand, it follows from \eqref{Butman-ee-cond} that
\[
    s_{i-k}=s_i\operatorname{sign}(x)^k,
\]
and hence
\[
    S_i(r) = s_i\sum_{k=1}^{p} b_k x^{-k} = s_i B(x).
\]
Consequently, the recursive sign-selection rule \eqref{Butman-sign-cond} is compatible with \eqref{Butman-ee-cond} precisely when
\begin{equation*}
    s_i=\operatorname{sign}(S_i(r)) = s_i\operatorname{sign}(B(x)),
\end{equation*}
which is equivalent to the compatibility condition
\begin{equation}\label{Butman-branch-2}
    B(x)\ge 0.
\end{equation}
Hence, the equation \eqref{butman-abs-eqn} is not valid without this compatibility condition \eqref{Butman-branch-2}. 

For $p=1$, this compatibility condition \eqref{Butman-branch-2} is automatic. Indeed,
$$
    B(x)=b_1x^{-1}.
$$
Given \(r>1\), choose $x=\operatorname{sign}(b_1)r$ when \(b_1\neq0\). Then
\[
    B(x)=\frac{|b_1|}{r}>0.
\]
Moreover,
\[
    S_i(r)=b_1s_{i-1}r^{-1},
\]
such that
\[
    \operatorname{sign}(S_i(r)) = \operatorname{sign}(b_1)s_{i-1} = \operatorname{sign}(x)s_{i-1} = s_i.
\]
Thus, \eqref{Butman-sign-cond} and \eqref{Butman-ee-cond} are consistent in the AR(1) case, and \eqref{Butman-one-step-factor} becomes
\[
    \left(1+|S_i(r)|\right)^2 = \left(1+\frac{|b_1|}{r}\right)^2,
\]
which is Butman's AR(1) recursion (\cite[Equation (4)]{butman1976linear}). 

For $p>1$, however, a real root $x$ of the absolute-value equation \eqref{butman-abs-eqn} need not satisfy the compatibility condition \eqref{Butman-branch-2}. Example~\ref{Butman-example-1:label} in Appendix~\ref{sec:example_of_ar_2_gaussian_channel} gives an AR(2) channel for which the absolute-value equation \eqref{butman-abs-eqn} has real roots but every such root has $B(x)<0$, so none is compatible. Thus, \eqref{butman-abs-eqn} is not, by itself, a generally valid achievable-rate characterization for AR($p$) Gaussian channels; an additional sign-compatibility condition is required. Separately, Example~\ref{Butman-example-2:label} in Appendix~\ref{sec:example_of_ar_2_gaussian_channel} gives an AR(2) channel for which $B(x)>0$ for every real $x\neq 0$, so the fixed replacement $|B(x)|=-B(x)$ leading to the polynomial equation \eqref{butman-final-eqn-1} is invalid. These examples show that neither the absolute-value equation \eqref{butman-abs-eqn} alone nor the fixed-sign polynomial equation \eqref{butman-final-eqn-1} is a branch-complete achievable-rate characterization for arbitrary AR($p$) Gaussian channels.

%Independently, even if one starts from \eqref{butman-abs-eqn}, the removal of the absolute value requires additional condition. Specifically, in our notation, Butman's expression \eqref{butman-final-eqn-1} replaced $|B(x)|$ by $-B(x)$, which is valid only at some root of \eqref{butman-final-eqn-1} and $B(x)\le0$. Example~\ref{Butman-example-2:label} in Appendix~\ref{sec:example_of_ar_2_gaussian_channel} shows that this fixed replacement is not valid in general. Consequently, without such condition, Butman's AR($p$) polynomial expression for $R^{(p)}_{ee}$ is not rigorously established achievable rate for arbitrary AR($p$) Gaussian channels.


%Butman further conjectured that $R^{(p)}_{ee}$ in \eqref{Butman-ARp-mutual-info} is tight. However, as shown above, the Butman's expression for AR($p$) case \eqref{butman-final-eqn-1} is not a rigorously justified for arbitrary AR($p$) Gaussian channels, because the derivation requires additional sign-compatibility and branch conditions. Therefore, before stating Butman's conjecture in a rigorous form, the lower bound itself must be corrected.

These observations and Example~\ref{Butman-example-1:label}--\ref{Butman-example-2:label} suggest that the AR($p$) part of Butman's derivation should return to \cite[Equation (39)]{butman1976linear}, before the absolute-value reduction is made. After the equal-energy specialization and the signed equal-ratio choice $x_i=x$, $|x|>1$, Butman's equation \cite[Equation (39)]{butman1976linear} becomes
\begin{align}
    x^2  &= 1+P\left(1+\sum_{k=1}^{p}b_k x^{-k}\right)^2 \nonumber\\
    &= 1+P L_Z^2(x^{-1}). \label{butman-corrected-signed-eqn}
\end{align}
Equivalently,
\begin{equation}\label{butman-corrected-power}
    P
    =
    \frac{x^2-1}{L_Z^2(x^{-1})}
    =
    \frac{x^2-1}
    {\left(1+\sum_{k=1}^{p}b_kx^{-k}\right)^2}.
\end{equation}
In particular, when \(x>1\), \eqref{butman-corrected-signed-eqn} reduces to the positive branch \eqref{butman-final-eqn-2} and the corresponding achievable rate is exactly $R^{(+)}_{\mathrm{Butman}}(P)$, whereas the case \(x<-1\) is obtained by evaluating \(L_Z(x^{-1})\) at a negative argument.

Thus, the corrected branch-complete achievable rate derived by Butman is
\begin{equation}\label{butman-branch-complete-rate}
R_{\mathrm{Butman}}(P) = \max_\gamma\log |\gamma|,
\end{equation}
where the maximum is taken over all $\gamma\in\mathbb{R}$ satisfying $|\gamma|>1$ and 
\begin{equation}\label{butman-branch-complete-power}
   (\gamma^2-1)H_Z^2(\gamma^{-1})=P,
\end{equation}
where \(H_Z(z)=1/L_Z(z)\). Comparing \eqref{butman-branch-complete-rate}--\eqref{butman-branch-complete-power} with Theorem~\ref{ISK1-thm}, we obtain
\begin{equation}\label{butman-sk1-identity}
    R_{\mathrm{Butman}}(P)=C_{\mathrm{SK1}}(P).
\end{equation}
Thus, the SK(1) coding scheme is the branch-complete and sign-consistent formulation of the constant signed-ratio correction to Butman's first-order equal-energy coding scheme. Its restriction to $\gamma>1$ gives the positive-branch expression for $R^{(+)}_{\mathrm{Butman}}(P)$ subject to \eqref{butman-final-eqn-2}, while allowing $\gamma<-1$ supplies the second real branch without removing an absolute value by a fixed sign.

Butman's original conjecture was that the AR($p$) equal-energy rate $R^{(p)}_{ee}$ in equation \eqref{Butman-ARp-mutual-info} equals the feedback capacity $C_\mathrm{FB}(P)$. Because $R^{(p)}_{ee}$ is not rigorously established as achievable for arbitrary AR($p$) Gaussian channels, we replace it by the rigorously achievable rate $R_{\mathrm{Butman}}(P)=C_{\mathrm{SK1}}(P)$ and formulate the following corrected form.

%In particular, Butman's positive-branch expression for $R^{(+)}_{\mathrm{Butman}}(P)$ subject to \eqref{butman-final-eqn-2} is only one branch result of the SK(1) coding scheme, while the full SK(1) formulation keeps both possible signs through the real parameter \(\gamma\) and hence avoids the ambiguity caused by dropping the absolute value.

Accordingly, the rigorous version of Butman's conjecture should be stated in terms of this corrected achievable rate $R_{\mathrm{Butman}}(P)$, or equivalently in terms of $C_{\mathrm{SK1}}(P)$.
\begin{conj}[Corrected form of Butman's conjecture]\label{Butman-conj:label}
For stationary AR($p$) Gaussian channels \eqref{channel-arp},
\[
    C_{\mathrm{FB}}(P)=C_{\mathrm{SK1}}(P).
\]
\end{conj}
For \(p=1\), this conjecture is true by the known feedback-capacity formula for ARMA(1,1) Gaussian channels \cite{kim2009feedback,yang2007feedback,sabag2023}. For \(p>1\), however, the conjecture is false in general. In particular, Kim \cite[Section IV]{Kim2006} observed, via an interleaving argument, that the Butman's conjecture cannot hold in general for $p>1$. Proposition~\ref{prop:interleaved-AR1} gives an explicit stationary AR(2) channel and a complete capacity calculation. In Section~\ref{sec:sk2-coding-scheme}, we also give another class of AR(2) examples for which the SK(2) coding scheme achieves a rate strictly larger than the rate achieved by the SK(1) coding scheme. Hence, for certain AR(2) Gaussian channels
\[
    C_{\mathrm{FB}}(P) \ge C_{\mathrm{SK2}}(P) > C_{\mathrm{SK1}}(P),
\]
which disproves Conjecture~\ref{Butman-conj:label}.

In the following section, we introduce the SK(2) coding scheme, in which the underlying message process follows a second-order deterministic recursion. This provides a strict generalization of the SK(1) coding scheme.


\section{The SK(2) Coding Scheme}\label{sec:sk2-coding-scheme}
Building on the framework of the SK(1) coding scheme, we now present its natural extension: the SK(2) coding scheme. This generalization is achieved by generalizing the underlying message process from a first-order to a second-order linear deterministic recursion. Formally, let $U_1$ and $U_2$ be two independent standard normal random variables, independent of the noise process $\{Z_i\}$. The transmitter initially sends 
$$
X_1 = U_1,\quad X_2 = U_2
$$
and subsequently sends 
\begin{equation}\label{SK2-coding-input}
X_i = V_i - \E{[V_i\mid Y_1^{i-1}]}, \quad i=3,4,\dots
\end{equation}
where $\{V_n\}$ is the message process satisfying 
\begin{equation}\label{SK2-coding-message}
\left\{ 
\begin{aligned}
V_{i+1} &= ~aV_i + bV_{i-1},\quad i\ge2,\\
V_2 ~~&= ~U_2,\\
V_1 ~~&= ~U_1,
\end{aligned}
\right.
\end{equation}
where $a,b\in\mathbb{R}$. Analogous to the definition of the SK(1) capacity in \eqref{feedback-capacity-SK1}, we define the {\em SK(2) information rate capacity}, or simply, {\em SK(2) capacity} as
\begin{equation*}
C_{\mathrm{SK2}}(P) = \sup_{(V,X)} \bar{I}(V;Y),
\end{equation*}
where the supremum is taken over all pairs $(V,X)$ satisfying (\ref{SK2-coding-input}), (\ref{SK2-coding-message}) and the constraint \eqref{power-constraint-sk1}. Thus, $C_{\mathrm{SK2}}(P)$ is the maximal achievable information rate
within the SK(2) coding class. In particular,
$$
    C_{\mathrm{SK2}}(P)\le C_{\mathrm{FB}}(P).
$$

Moreover, when $b=0$, the recursion \eqref{SK2-coding-message} reduces to $V_{n+1}=aV_n$, essentially recovering the SK(1) coding scheme \eqref{SK1-coding-message}. Hence, the SK(1) coding scheme is contained in the SK(2) coding class as a special case, and consequently
\[
    C_{\mathrm{SK1}}(P)\le C_{\mathrm{SK2}}(P)\le C_{\mathrm{FB}}(P).
\]


Thus, the case of interest is $b\neq 0$. To proceed, we define the {\em genuine SK(2) capacity} under the average power constraint $P$ as
\begin{equation*}
C^{\mathrm{g}}_{\mathrm{SK2}}(P)=\sup_{(V,X)}\bar{I}(V;Y),
\end{equation*}
where the supremum is taken over all pairs $(V,X)$ satisfying \eqref{SK2-coding-input}, \eqref{SK2-coding-message}, $b\neq 0$, and the constraint \eqref{power-constraint-sk1}.


As with the SK(1) coding scheme, the parameters $a$ and $b$ must satisfy certain constraints. Rather than imposing conditions directly on $a$ and $b$, we analyze the corresponding characteristic equation (in $\lambda$) of \eqref{SK2-coding-message}:
\begin{equation}\label{sk2-char-eqn}
   \lambda^2 - a\lambda - b =0,
\end{equation}
and impose the conditions on its roots, $\gamma_1$ and $\gamma_2$. Then $\{V_n\}$ admits the finite-rank representation
\begin{equation*}\label{Vn-genaral-sol}
   V_n = a_n U_1 +b_n U_2,\quad n\ge1,
\end{equation*}
where the deterministic coefficients $(a_n,b_n)$ are given as follows, according to the nature of the roots.

\textbf{Case A:} $\gamma_1 \neq \gamma_2$ (distinct roots). Then 
\begin{equation}\label{Vn-case-a}
   a_n = -(c_1\gamma_2\gamma_1^n +c_2\gamma_1\gamma_2^n),\quad b_n = c_1\gamma_1^n +c_2\gamma_2^n,
\end{equation}
where
$$
   c_1 = \frac{1}{\gamma_1(\gamma_1-\gamma_2)},\quad c_2 = \frac{1}{\gamma_2(\gamma_2-\gamma_1)}.
$$

\textbf{Case B:} $\gamma_1 = \gamma_2 =: \gamma$ (repeated root). Then
\begin{equation}\label{Vn-case-b}
   a_n = (2-n)\gamma^{n-1},\quad b_n = (n-1)\gamma^{n-2}.
\end{equation}

% \begin{enumerate}
%     \item[] \textbf{Case A:} \(\gamma_1 \neq \gamma_2\) (distinct roots);
    
%     \item[] \textbf{Case B:} \(\gamma_1 = \gamma_2 =: \gamma\) (repeated root).
% \end{enumerate}
% In either case, $\{V_n\}$ admits the finite-rank representation
% \begin{equation*}\label{Vn-genaral-sol}
%    V_n = a_n U_1 +b_n U_2,\quad n\ge1,
% \end{equation*}
% where the deterministic coefficients $(a_n,b_n)$ are given by:
% \begin{itemize}
% \item[] \textbf{Case A:}
% \begin{equation}\label{Vn-case-a}
%    a_n = -(c_1\gamma_2\gamma_1^n +c_2\gamma_1\gamma_2^n),\quad b_n = c_1\gamma_1^n +c_2\gamma_2^n,
% \end{equation}
% with $c_1,c_2$ defined as
% \[
%    c_1 = \frac{1}{\gamma_1(\gamma_1-\gamma_2)},\quad c_2 = \frac{1}{\gamma_2(\gamma_2-\gamma_1)}.
% \]
% \item[] \textbf{Case B:} 
% \begin{equation}\label{Vn-case-b}
%    a_n = (2-n)\gamma^{n-1},\quad b_n = (n-1)\gamma^{n-2}.
% \end{equation}
% \end{itemize} 
Clearly, $a$ and $b$ are uniquely determined by $\gamma_1$ and $\gamma_2$; hence, it suffices to constrain the pair $(\gamma_1,\gamma_2)$. As shown in the following theorem, it is natural to require that $|\gamma_1|\wedge|\gamma_2| >1$, where we have adopted the notation $\phi\wedge \psi\triangleq \min(\phi,\psi)$.

The following theorem provides a closed-form characterization of the genuine SK(2) capacity $C^\mathrm{g}_\mathrm{SK2}(P)$, which yields a lower bound on $C_\mathrm{SK2}(P)$.

%We now state the corresponding achievable rate $\bar{I}_{\mathrm{SK2}}(P)$ for the SK(2) coding scheme under the condition $b\neq 0$ in \eqref{SK2-coding-message} over stationary AR$(p)$ Gaussian channels, yielding a lower bound on $C_\mathrm{SK2}(P)$.
%Next, we shall provide the result regarding the achievable rates achieved by the SK(2) coding scheme when $b\neq 0$ in \eqref{SK2-coding-message}.
\begin{thm}\label{main-theorem-sk2}
Suppose the noise $\{ Z_i \}$ is an $AR(p)$ Gaussian process with parameters $\beta_k$, $\vert \beta_k\vert<1$ for all $k=1,2,...,p$, namely, it has the power spectral density
\begin{equation*}
S_Z(e^{i\theta}) = \vert H_Z(e^{i\theta}) \vert^2 = \frac{1}{\vert L_Z(e^{i\theta}) \vert^2} = \frac{1}{\vert\prod_{k=1}^{p}(1+\beta_ke^{i\theta})\vert^2}.
\end{equation*}
Then, we have
\begin{equation}\label{I-SK2-thm}
C^\mathrm{g}_{\mathrm{SK2}}(P) =\sup_{\gamma_1,\gamma_2} \log  |\gamma_1\gamma_2|
\end{equation}
and the supremum is taken over all $\gamma_1,\gamma_2\in\mathbb{C}$ satisfying $|\gamma_1|\wedge|\gamma_2|>1$, $\gamma_1+\gamma_2\in \mathbb{R}$, $\gamma_1\gamma_2\in\mathbb{R}$, and the power constraint 
\begin{equation}\label{power-eqn-sk2-thm}
\frac{1}{\Delta}\left( \frac{L_{11}}{\gamma^2_2}+\frac{L_{22}}{\gamma_1^2} +\frac{2L_{12}}{\gamma_1\gamma_2}  \right)= P,
\end{equation}
where 
\begin{equation}\label{main-thm-notations}
\begin{aligned}
\Delta &= (\gamma_1-\gamma_2)^2(L_{11}L_{22}-L_{12}^2),\\
L_{ij} &=\frac{c_{i}c_{j}L_Z(\gamma_i^{-1})L_Z(\gamma_j^{-1})}{\gamma_i\gamma_j-1} \quad \mbox{for $i,j=1,2,$}\\
c_1 &= \frac{1}{\gamma_1(\gamma_1-\gamma_2)}, \quad c_2 = \frac{1}{\gamma_2(\gamma_2-\gamma_1)}.
\end{aligned}
\end{equation}
\end{thm}

\begin{proof}
Consider the SK(2) coding scheme formulated by \eqref{SK2-coding-input} and \eqref{SK2-coding-message} over the stationary AR$(p)$ Gaussian channel $Y_i = X_i +Z_i$ with the noise spectral density function $S_Z(e^{i\theta})$ given by (\ref{sdf-arp}).

We first address Case A. Recall that 
\begin{equation*}
\begin{aligned}
Y_n &= X_n + Z_n= V_n -\E[V_n\mid Y_1^{n-1}] +Z_n
\end{aligned}
\end{equation*}
and define
\begin{equation*}
Y_n^\ast = V_n + Z_n.
\end{equation*}
Clearly, $Y_n^\ast = Y_n +\E[V_n\mid Y_1^{n-1}]$ and therefore, for any $i\ge 1$, $\{ Y_1^i \}$ and $\{Y_1^{\ast,i}\}$ are uniquely determined by each other. Furthermore, define
\begin{equation}\label{Y-tlta}
\widetilde{Y}_n =\left\{ 
\begin{aligned}
& Y_n^\ast, ~~~~~~~~~~ n\le p;\\[2pt]
&L(B)Y^\ast_n, ~~~n\ge p+1,
\end{aligned}
\right.
\end{equation}
where $L$ stands for $L_Z$, $B$ denotes the backward-shift operator. Now by (\ref{Vn-case-a}), we can express 
\begin{equation}\label{Y-tlta-2}
\begin{aligned}
\widetilde{Y}_n & = L(B)(V_n +Z_n)\\
& =  L(B)(a_n U_1 +b_nU_2) +L(B)Z_n\\
&= d_{1,n}U_1 +d_{2,n}U_2 +W_n
\end{aligned}
\end{equation}
for $n\ge p+1$, where
\begin{equation}\label{d1d2}
\left\{ 
\begin{aligned}
d_{1,n} &=-c_1\gamma_2\gamma_1^nL(\gamma_1^{-1})- c_2\gamma_1\gamma_2^n L(\gamma_2^{-1});\\[2pt]
d_{2,n} &= c_1\gamma_1^nL(\gamma_1^{-1}) + c_2\gamma_2^n L(\gamma_2^{-1}).
\end{aligned}
\right.
\end{equation}
Thus, by \eqref{Y-tlta}, \eqref{Y-tlta-2} and \eqref{d1d2}, we have 
\begin{equation}\label{Y-tlta-re}
\widetilde{Y}_n =\left\{ 
\begin{aligned}
& a_n U_1 +b_n U_2 + Z_n, ~~~~~~~~~~~~ n\le p;\\[2pt]
&d_{1,n}U_1 + d_{2,n}U_2 + W_n,~~~~ ~~~n\ge p+1.
\end{aligned}
\right.
\end{equation}

For $n\ge p+1$, we have the following chain of equalities:
%The mutual information rate between $V_1^n$ and $Y_1^n$ for $n\ge p+1$ is given by
\begin{equation}\label{mutual-info-sk2}
\begin{aligned}
\frac{1}{n}I(V_1^n;Y_1^n) &\overset{(a)}{=}\frac{1}{n}I(U_1,U_2;Y_1^n)\\
&\overset{(b)}{=}\frac{1}{n}I(U_1,U_2;\widetilde{Y}_1^n)      \\
&= \frac{1}{n}I(U_1;\widetilde{Y}_1^n) +\frac{1}{n}I(U_2;\widetilde{Y}_1^n\mid U_1)\\
&\overset{(c)}{=}-\frac{1}{2n}\log\E[\vert U_1-\E[U_1\mid \widetilde{Y}_{1}^n] \vert^2] -\frac{1}{2n}\log\E[\vert U_2-\E[U_2\mid U_1, \widetilde{Y}_{1}^n] \vert^2],
\end{aligned}
\end{equation}
where $(a)$ follows from \eqref{SK2-coding-message}, $(b)$ follows from the fact that $Y_1^n$ and $\widetilde{Y}_1^n$ are uniquely determined by each other and $(c)$ follows from the relation between mutual information and MMSE for Gaussian random variables, together with the entropy formula for a multivariate normal distribution (see, e.g., \cite[Theorem 8.6.6 \& its corollary]{cover1991elements}).

Next, for $n\ge p+1$, we also have
\begin{equation}\label{power-eqn-sk2}
   \begin{aligned}
      \E[|X_n|^2] &= \E[\vert V_n-\E[V_n\mid Y_1^{n-1}]\vert ^2]\\
      &\overset{(d)}{=}\E[\vert V_n - \E[ V_n\mid \widetilde{Y}_1^{n-1}]\vert^2],
   \end{aligned}
\end{equation}
where $(d)$ follows from the fact that $Y_1^n$ and $\widetilde{Y}_1^n$ are uniquely determined by each other.

Let 
\[
A = \begin{bmatrix} a_1 & a_2 & \cdots & a_p \end{bmatrix}^T,\qquad
B = \begin{bmatrix} b_1 & b_2 & \cdots & b_p \end{bmatrix}^T,
\]
and
\[
D_{1,n} = \begin{bmatrix} d_{1,p+1} & d_{1,p+2} & \cdots & d_{1,n} \end{bmatrix}^T,\qquad
D_{2,n} = \begin{bmatrix} d_{2,p+1} & d_{2,p+2} & \cdots & d_{2,n} \end{bmatrix}^T.
\]
Since $U_1$ and $\{\widetilde{Y}_1^n\}$ are jointly Gaussian, the conditional mean squared error is given by 
\begin{equation}\label{MMSE-U-1-a}
\begin{aligned}
   \E[\vert U_1-\E[U_1\mid \widetilde{Y}_1^n] \vert^2] &= \mathrm{Var}(U_1) - \mathrm{Cov}(U_1,\widetilde{Y}_1^n)(\mathrm{Var}(\widetilde{Y}_1^n))^{-1}\mathrm{Cov}(\widetilde{Y}_1^n,U_1)\\
   &= 1 - \begin{bmatrix}A^T & D_{1,n}^T \end{bmatrix}\Sigma_{\widetilde{Y}_1^n}^{-1}\begin{bmatrix}A\\[2pt] D_{1,n}\end{bmatrix},
\end{aligned}
\end{equation}
where the second equality uses $\mathrm{Var}(U_1)=1$ and the fact that
\[
   \mathrm{Cov}(U_1,\widetilde{Y}_1^n)= \begin{bmatrix}A^T & D_{1,n}^T \end{bmatrix}.
\]
Here $\Sigma_{\widetilde{Y}_1^n} \triangleq\mathrm{Var}(\widetilde{Y}_1^n)$ is given explicitly by
\begin{equation}\label{def-Sigma-Y-tilta}
   \Sigma_{\widetilde{Y}_1^n} =\begin{bmatrix}
         AA^T+BB^T+\Sigma_Z & A D_{1,n}^T+BD_{2,n}^T\\[2pt]
         D_{1,n}A^T +D_{2,n} B^T & I_{n-p}+D_{1,n}D_{1,n}^T +D_{2,n}D_{2,n}^T
      \end{bmatrix},
\end{equation}
where $\Sigma_Z\triangleq\mathrm{Var}(Z_1^p)$, and $I_{n-p}$ denotes the $(n-p)\times(n-p)$ identity matrix. 

Moreover, let 
\[
   \widetilde{D}_{1,n}=\begin{bmatrix}A \\[2pt]D_{1,n}\end{bmatrix}, \qquad\widetilde{D}_{2,n}=\begin{bmatrix}B\\[2pt]D_{2,n}\end{bmatrix},\qquad\widetilde{\Sigma}_{Z,n}=\begin{bmatrix}\Sigma_Z&0\\[2pt]0 &I_{n-p}\end{bmatrix}.
\]
Then, the equation \eqref{MMSE-U-1-a} can be rewritten in the following compact form: 
\begin{equation}\label{MMSE-U-1}
\begin{aligned}
   \E[\vert U_1-\E[U_1\mid \widetilde{Y}_1^n] \vert^2] &=1 - \widetilde{D}_{1,n}^T\left(\widetilde{\Sigma}_{Z,n} +\widetilde{D}_{1,n}\widetilde{D}_{1,n}^T +\widetilde{D}_{2,n}\widetilde{D}_{2,n}^T\right)^{-1}\widetilde{D}_{1,n}.
\end{aligned}
\end{equation}
Note that by \cite[Proposition 5.1.1]{brockwell2009time}, $\Sigma_Z$ is non-singular and hence positive definite. Then, applying Lemma \ref{lemma1} to \eqref{MMSE-U-1} yields 
\begin{equation}\label{MMSE-U-2}
\begin{aligned}
   \E[\vert U_1-\E[U_1\mid \widetilde{Y}_1^n]\vert^2]&=\frac{1+\widetilde{D}_{2,n}^T\widetilde{\Sigma}_{Z,n}^{-1}\widetilde{D}_{2,n}}{\det\begin{pmatrix}
      1+\widetilde{D}^T_{1,n}\widetilde{\Sigma}_{Z,n}^{-1}\widetilde{D}_{1,n} & \widetilde{D}^T_{1,n}\widetilde{\Sigma}_{Z,n}^{-1}\widetilde{D}_{2,n}\\[2pt]
      \widetilde{D}^T_{2,n}\widetilde{\Sigma}_{Z,n}^{-1}\widetilde{D}_{1,n} & 1+ \widetilde{D}^T_{2,n} \widetilde{\Sigma}_{Z,n}^{-1}\widetilde{D}_{2,n}
   \end{pmatrix}}\\
   &=\frac{1+ B^T\Sigma_Z^{-1}B +D_{2,n}^TD_{2,n}}{\det\begin{pmatrix}1+A^T\Sigma_Z^{-1}A +D^T_{1,n}D_{1,n} & A^T\Sigma_Z^{-1}B+D^T_{1,n}D_{2,n}\\[2pt]
      B^T\Sigma_Z^{-1}A+D^T_{2,n}D_{1,n} &1+B^T\Sigma_Z^{-1}B+D^T_{2,n}D_{2,n}
   \end{pmatrix}}.
\end{aligned}
\end{equation}
% \begin{equation}\label{MMSE-U-2}
% \E[\vert U_1-\E[U_1\vert \widetilde{Y}_1^n]\vert^2] = \frac{1+D_{2,n}^TD_{2,n}}{\vert I_2 +D_n^TD_n\vert},
% \end{equation}
%where $I_2$ is the $2\times 2$ identity matrix and $D_n=[D_{1,n} ~D_{2,n}]$.
Similarly, for $U_2$ we obtain
\begin{equation}\label{MMSE-V-2}
\begin{aligned}
\E[\vert U_2-\E[U_2\mid U_1, \widetilde{Y}_1^n]\vert^2] &= 1 - \begin{bmatrix}0 & \widetilde{D}^T_{2,n}\end{bmatrix} \begin{bmatrix}1 & \widetilde{D}_{1,n}^T\\[2pt] \widetilde{D}_{1,n} & \widetilde{\Sigma}_{Z,n}+\widetilde{D}_{1,n}\widetilde{D}_{1,n}^T+\widetilde{D}_{2,n}\widetilde{D}_{2,n}^T\end{bmatrix}^{-1} \begin{bmatrix}
   0\\ \widetilde{D}_{2,n}
\end{bmatrix}        \\
&= 1 -   \widetilde{D}_{2,n}^T(\widetilde{\Sigma}_{Z,n}+\widetilde{D}_{2,n}\widetilde{D}_{2,n}^T)^{-1}\widetilde{D}_{2,n}         \\
&=\frac{1}{1+\widetilde{D}_{2,n}^T\widetilde{\Sigma}_{Z,n}^{-1}\widetilde{D}_{2,n}}\\
&=\frac{1}{1+B^T\Sigma_Z^{-1}B+D^T_{2,n}D_{2,n}}.
\end{aligned}
\end{equation}
Consequently, it follows from \eqref{MMSE-U-2} and \eqref{MMSE-V-2} that \eqref{mutual-info-sk2} can be simplified as
\begin{equation}\label{pf-sk2-mutual-info}
   \frac{1}{n}I(V_1^n;Y_1^n)= \frac{1}{2n}\log \det(\Delta_{2,n}),
\end{equation}
where $\Delta_{2,n}=I_2 +\widetilde{D}_n^T\widetilde{\Sigma}_{Z,n}^{-1}\widetilde{D}_n$ and $\widetilde{D}_n =\begin{bmatrix}\widetilde{D}_{1,n} & \widetilde{D}_{2,n}\end{bmatrix}$.

On the other hand, since $V_n$ and $\{\widetilde{Y}_1^{n-1}\}$ are also jointly Gaussian, it follows from \cite[Theorem 4.6]{gray2004introduction} that
\begin{equation}\label{MMSE-Vn}
\begin{aligned}
\E[| V_n-\E[V_n\mid \widetilde{Y}_1^{n-1}]|^2]&= \mathrm{Var}(V_n) - \mathrm{Cov}(V_n,\widetilde{Y}_1^{n-1})\left(\mathrm{Var}(\widetilde{Y}_1^{n-1})\right)^{-1}\mathrm{Cov}(\widetilde{Y}_1^{n-1},V_n) \\
&= a_n^2 +b_n^2 -\begin{bmatrix}a_n&b_n\end{bmatrix}\widetilde{D}^T_{n-1}\Sigma_{\widetilde{Y}_1^{n-1}}^{-1}\widetilde{D}_{n-1}\begin{bmatrix}a_n\\[2pt]b_n\end{bmatrix},
\end{aligned}
\end{equation}
where the second equality uses $\mathrm{Var}(V_n)=a_n^2+b_n^2$ and the fact that
\begin{equation*}
\begin{aligned}
   \mathrm{Cov}(V_n,\widetilde{Y}_1^{n-1}) &= \begin{bmatrix}a_n A^T+b_nB^T& a_nD_{1,n-1}^T+b_nD_{2,n-1}^T\end{bmatrix}\\
   &=\begin{bmatrix}a_n &b_n \end{bmatrix}\begin{bmatrix}\widetilde{D}^T_{1,n-1}\\[2pt]\widetilde{D}^T_{2,n-1}\end{bmatrix}\\
   &=\begin{bmatrix}a_n&b_n\end{bmatrix}\widetilde{D}^T_{n-1}.
\end{aligned}
\end{equation*}
Now, by \eqref{def-Sigma-Y-tilta} with $n$ replaced by $n-1$, the covariance matrix of $\widetilde{Y}_1^{n-1}$ can be expressed as
\begin{equation*}
\begin{aligned}
   \Sigma_{\widetilde{Y}_1^{n-1}}&=\widetilde{\Sigma}_{Z,n-1}+\widetilde{D}_{1,n-1}\widetilde{D}^T_{1,n-1}+\widetilde{D}_{2,n-1}\widetilde{D}^T_{2,n-1}\\
&=\widetilde{\Sigma}_{Z,n-1}+\widetilde{D}_{n-1}\widetilde{D}^T_{n-1},
\end{aligned}
\end{equation*}
which, together with \eqref{MMSE-Vn}, yields
\begin{equation}\label{MMSE-Vn-final}
   \begin{aligned}
      \E[\vert& V_n-\E[V_n\mid \widetilde{Y}_1^{n-1}]\vert^2]=a_n^2+b_n^2-\begin{bmatrix}a_n&b_n\end{bmatrix}\widetilde{D}^T_{n-1}\left(\widetilde{\Sigma}_{Z,n-1}+\widetilde{D}_{n-1}\widetilde{D}^T_{n-1}\right)^{-1}\widetilde{D}_{n-1}\begin{bmatrix}a_n\\[2pt]b_n\end{bmatrix}\\ 
&\overset{(e)}{=}a_n^2+b_n^2-\begin{bmatrix}a_n &b_n\end{bmatrix}\left(I_2 - (I_2 +\widetilde{D}^T_{n-1}\widetilde{\Sigma}_{Z,n-1}^{-1}\widetilde{D}_{n-1})^{-1}\right)\begin{bmatrix}a_n\\[2pt]b_n\end{bmatrix}\\
&=\begin{bmatrix}a_n&b_n\end{bmatrix}\left(I_2 +\widetilde{D}^T_{n-1}\widetilde{\Sigma}_{Z,n-1}^{-1}\widetilde{D}_{n-1}\right)^{-1}      \begin{bmatrix}a_n\\[2pt]b_n\end{bmatrix}\\
&=\begin{bmatrix}a_n&b_n\end{bmatrix}(\Delta_{2,n-1})^{-1} \begin{bmatrix}a_n\\[2pt]b_n\end{bmatrix},
\end{aligned}
\end{equation}
where $(e)$ follows from Lemma \ref{lemma1}. Now, note that
\begin{equation*}
\begin{aligned}
\Delta_{2,n-1}&=\begin{bmatrix}1&0\\[2pt]0&1\end{bmatrix}+ \begin{bmatrix}A^T&D_{1,n-1}^T\\[2pt]B^T&D_{2,n-1}^T \end{bmatrix}
\begin{bmatrix}\Sigma_Z^{-1}&0\\[2pt]0&I_{n-p-1}\end{bmatrix}\begin{bmatrix}A&B\\[2pt]D_{1,n-1}&D_{2,n-1}\end{bmatrix}\\
&=\begin{bmatrix}
1+A^T\Sigma_Z^{-1}A +D_{1,n-1}^TD_{1,n-1} & A^T\Sigma_Z^{-1}B+D_{1,n-1}^TD_{2,n-1}\\[2pt]
B^T\Sigma_Z^{-1}A +D^T_{2,n-1}D_{1,n-1}   & 1+B^T\Sigma_Z^{-1}B+D_{2,n-1}^TD_{2,n-1}
\end{bmatrix}.
\end{aligned}
\end{equation*}
Its inverse can be written as
\begin{equation*}
\hspace{-0.8cm}
\begin{aligned}
 (\Delta_{2,n-1})^{-1}&=\frac{1}{\det(\Delta_{2,n-1})}\begin{bmatrix}
    1+B^T\Sigma_Z^{-1}B+D_{2,n-1}^TD_{2,n-1}& -A^T\Sigma_Z^{-1}B-D_{1,n-1}^TD_{2,n-1}\\
    -B^T\Sigma_Z^{-1}A -D^T_{2,n-1}D_{1,n-1}&1+A^T\Sigma_Z^{-1}A +D_{1,n-1}^TD_{1,n-1}
 \end{bmatrix}\\
 &=\frac{1}{\det(\Delta_{2,n-1})}\left(\begin{bmatrix}1+D_{2,n-1}^TD_{2,n-1} & -D^T_{1,n-1}D_{2,n-1}\\ -D^T_{2,n-1}D_{1,n-1} & 1+D_{1,n-1}^TD_{1,n-1}\end{bmatrix} +\begin{bmatrix}B^T\Sigma_Z^{-1}B & -A^T\Sigma_Z^{-1}B \\ -B^T\Sigma_Z^{-1}A & A^T\Sigma_Z^{-1}A \end{bmatrix}              \right),
\end{aligned}
\end{equation*}
which, substituted into \eqref{MMSE-Vn-final} and then combined with \eqref{power-eqn-sk2}, gives
\begin{equation}\label{MMSE-Vn2}
\begin{aligned}
 \E[|X_n|^2]&= \frac{a_n^2+b_n^2+\|a_nD_{2,n-1}-b_nD_{1,n-1} \|^2+(a_nB^T-b_nA^T)\Sigma_Z^{-1}(a_nB-b_nA)}{\det(\Delta_{2,n-1})},
\end{aligned}
\end{equation}
where $\|\cdot\|$ denotes the Euclidean norm.

We are now ready to evaluate the limits of \eqref{pf-sk2-mutual-info} and (\ref{MMSE-Vn2}). Since $d_{1,n}$ and $d_{2,n}$ can be alternatively expressed as 
\begin{equation*}
\begin{aligned}
d_{1,n}=\begin{bmatrix}-\gamma_2 &-\gamma_1\end{bmatrix} \begin{bmatrix} c_1L(\gamma_1^{-1})\gamma_1^n\\[2pt]c_2L(\gamma_2^{-1})\gamma_2^n\end{bmatrix}, ~
d_{2,n}=\begin{bmatrix}1&1 \end{bmatrix}\begin{bmatrix} c_1L(\gamma_1^{-1})\gamma_1^n\\[2pt]c_2L(\gamma_2^{-1})\gamma_2^n\end{bmatrix},
\end{aligned}
\end{equation*}
it follows that
\begin{equation}\label{2norm-D1n}
\begin{aligned}
\|D_{1,n}\|^2&=D^T_{1,n}D_{1,n}     \\
 &=\sum_{k=p+1}^n\vert d_{1,k}\vert^2\\
&=\sum_{k=p+1}^n\begin{bmatrix}-\gamma_2&-\gamma_1\end{bmatrix}\begin{bmatrix}c_1L(\gamma_1^{-1})\gamma_1^k\\[2pt]c_2L(\gamma_2^{-1})\gamma_2^k\end{bmatrix}\begin{bmatrix}c_1L(\gamma_1^{-1})\gamma_1^k&c_2L(\gamma_2^{-1})\gamma_2^k\end{bmatrix}\begin{bmatrix}-\gamma_2\\[2pt]-\gamma_1\end{bmatrix} \\[2pt]
&=\sum_{k=p+1}^n \begin{bmatrix}\gamma_2&\gamma_1\end{bmatrix}\begin{bmatrix} c_1^2L^2(\gamma_1^{-1})\gamma_1^{2k} & c_1c_2L(\gamma_1^{-1})L(\gamma_2^{-1})\gamma_1^k \gamma_2^k\\[2pt]c_1c_2L(\gamma_1^{-1})L(\gamma_2^{-1})\gamma_1^k \gamma_2^k & c_2^2L^2(\gamma_2^{-1})\gamma_2^{2k}\end{bmatrix} \begin{bmatrix}\gamma_2\\[2pt]\gamma_1\end{bmatrix}\\[2pt]
&=\begin{bmatrix}\gamma_2&\gamma_1\end{bmatrix} 
\begin{bmatrix}
L_{11}\gamma_1^2(\gamma_1^{2n}-1) &L_{12}\gamma_1\gamma_2(\gamma_1^n\gamma_2^n-1)\\[2pt]L_{12}\gamma_1\gamma_2(\gamma_1^n\gamma_2^n-1) &L_{22}\gamma_2^2(\gamma_2^{2n}-1)
\end{bmatrix}
\begin{bmatrix}\gamma_2\\[2pt]\gamma_1\end{bmatrix}-R_1\\[2pt]
&=\begin{bmatrix}\gamma_2&\gamma_1\end{bmatrix}L_{n}^{(2)}\begin{bmatrix}\gamma_2\\[2pt]\gamma_1\end{bmatrix}-R_1,
\end{aligned}
\end{equation}
where we have defined $R_1=\sum_{k=1}^p |d_{1,k}|^2$ and
$$
L_{n}^{(2)}=\begin{bmatrix}
L_{11}\gamma_1^2(\gamma_1^{2n}-1) &L_{12}\gamma_1\gamma_2(\gamma_1^n\gamma_2^n-1)\\[2pt]L_{12}\gamma_1\gamma_2(\gamma_1^n\gamma_2^n-1) &L_{22}\gamma_2^2(\gamma_2^{2n}-1)
\end{bmatrix}.
$$
Next, set $R_2 = \sum_{k=1}^p |d_{2,k}|^2$ and $R_{12} = \sum_{k=1}^p d_{1,k} d_{2,k}$. Then, in complete analogy to the derivations for $D^T_{1,n}D_{1,n}$, we obtain
\begin{equation}\label{2norm-D2n}
\begin{aligned}
D_{2,n}^TD_{2,n}=\begin{bmatrix}1&1\end{bmatrix} L_n^{(2)}\begin{bmatrix}1\\[2pt]1\end{bmatrix}-R_2,~
D_{1,n}^TD_{2,n}=\begin{bmatrix}-\gamma_2&\gamma_1\end{bmatrix} L_n^{(2)}\begin{bmatrix}1\\[2pt]1\end{bmatrix}-R_{12}.
\end{aligned}
\end{equation}
Hence, 
\begin{equation}\label{diff-D12n}
\begin{aligned}
   D_{1,n}^T&D_{1,n}D_{2,n}^TD_{2,n}-D_{1,n}^TD_{2,n}D_{2,n}^TD_{1,n}=\sum_{p+1\le i<j}(d_{1,i}d_{2,j}-d_{1,j}d_{2,i})^2\\
   &=(\gamma_1-\gamma_2)^2c_1^2c_2^2\gamma_1^2\gamma_2^2L^2(\gamma_1^{-1})L^2(\gamma_2^{-1})\left(\frac{(\gamma_1^{2n}-1)(\gamma_2^{2n}-1)}{(\gamma_1^2-1)(\gamma_2^2-1)} - \frac{(\gamma_1^n\gamma_2^n-1)^2}{(\gamma_1\gamma_2-1)^2}  \right)\\[2pt]
   &~~~~-R_1 D^T_{2,n}D_{2,n}-R_2 D^T_{1,n}D_{1,n}+2R_{12}D^T_{1,n}D_{2,n}-R_1R_2+R^2_{12}\\
   &=\Delta \gamma_1^2\gamma_2^2(\gamma_1^{2n}\gamma_2^{2n}-1) + 2L^2_{12}\gamma_1^2\gamma_2^2(\gamma_1-\gamma_2)^2\gamma_1^n\gamma_2^n -\frac{(\gamma_1\gamma_2-1)^2\Delta \gamma_1^2\gamma_2^2}{(\gamma_1-\gamma_2)^2}(\gamma_1^{2n}+\gamma_2^{2n})\\[2pt]
   &~~~~-R_1 D^T_{2,n}D_{2,n}-R_2 D^T_{1,n}D_{1,n}+2R_{12}D^T_{1,n}D_{2,n}-R_1R_2+R^2_{12}.
\end{aligned}
\end{equation}
Furthermore, we have
\begin{equation}\label{det-sk2-1}
\hspace{-0.5cm}
\begin{aligned}
\det(\Delta_{2,n}) &= \det(I_2+\widetilde{D}_n^T\widetilde{\Sigma}_{Z,n}^{-1}\widetilde{D}_n)       \\
&= \det \begin{pmatrix}
   1+D_{1,n}^TD_{1,n} +A^T\Sigma_Z^{-1}A & D_{1,n}^TD_{2,n}+A^T\Sigma_Z^{-1}B\\[2pt] D_{2,n}^TD_{1,n} +B^T\Sigma_Z^{-1}A & 1+D_{2,n}^TD_{2,n}+B^T\Sigma_Z^{-1}B \end{pmatrix}\\
   &=(1+D_{1,n}^TD_{1,n})(1+D_{2,n}^TD_{2,n}) -D_{1,n}^TD_{2,n}D_{2,n}^TD_{1,n} +B^T\Sigma_Z^{-1}B(1+D^T_{1,n}D_{1,n}) \\[2pt]
   &~~~~+A^T\Sigma_Z^{-1}A(1+D^T_{2,n}D_{2,n}) +A^T\Sigma_Z^{-1}AB^T\Sigma_Z^{-1}B - 2A^T\Sigma_Z^{-1}BD^T_{1,n}D_{2,n}-(A^T\Sigma_Z^{-1}B)^2,
    \end{aligned}
\end{equation}
which, together with (\ref{2norm-D1n}) (\ref{2norm-D2n}) and (\ref{diff-D12n}), immediately implies that
\begin{equation}\label{det-sk2}
\det(\Delta_{2,n})= \Delta \gamma_1^2\gamma_2^2\gamma_1^{2n}\gamma_2^{2n} +O\left((|\gamma_1|\vee|\gamma_2|)^{2n}\right),
\end{equation}
where we have adopted the notation $\phi\vee \psi\triangleq \max(\phi,\psi)$.

Consequently, it follows from (\ref{det-sk2}) that, upon taking the limit as $n\to\infty$ in \eqref{pf-sk2-mutual-info}, the mutual information rate achieved by the SK(2) coding scheme specified in equations (\ref{SK2-coding-input}) and (\ref{SK2-coding-message}) is given by
\begin{equation}\label{mutual-info-sk2-proof}
\begin{aligned}
            \bar{I}(V;Y)&=\lim_{n\to\infty}\frac{1}{n}I(V_1^n;Y_1^n)\\
            &= \log|\gamma_1 \gamma_2|.
\end{aligned}
\end{equation}

On the other hand, we now examine the limiting behavior of \eqref{MMSE-Vn2} as \(n \to \infty\). To this end, observe that
$$
a_nd_{2,k}-b_nd_{1,k}=\begin{bmatrix}\gamma_2^{n-1}&\gamma_1^{n-1}\end{bmatrix} \begin{bmatrix}
   c_1L(\gamma_1^{-1})\gamma_1^k \\[2pt]c_2L(\gamma_2^{-1})\gamma_2^k
\end{bmatrix},
$$
which implies that 
\begin{equation}\label{aD-bD}
\begin{aligned}
\|&a_nD_{2,n-1}-b_nD_{1,n-1} \|^2= \begin{bmatrix}
      \gamma_2^{n-1}&\gamma_1^{n-1}
\end{bmatrix}
L_{n-1}^{(2)}
\begin{bmatrix}
      \gamma_2^{n-1}\\[2pt]\gamma_1^{n-1}
\end{bmatrix} - \begin{bmatrix}\gamma_2^{n-1}&\gamma_1^{n-1}\end{bmatrix}L_p^{(2)}\begin{bmatrix}\gamma_2^{n-1}\\[2pt]\gamma_1^{n-1}\end{bmatrix}\\
&=L_{11}\gamma_1^2\gamma_2^{2(n-1)}(\gamma_1^{2(n-1)}-1) +L_{22}\gamma_2^2\gamma_1^{2(n-1)}(\gamma_2^{2(n-1)}-1)+2L_{12}\gamma_1^{n}\gamma_2^{n}(\gamma_1^{n-1}\gamma_2^{n-1}-1) \\[2pt]
&~~~~-(L_p^{(2)})_{11}\gamma_2^{2(n-1)}-(L_p^{(2)})_{22}\gamma_1^{2(n-1)}-2(L_p^{(2)})_{12}\gamma^{n-1}_1\gamma^{n-1}_2.
\end{aligned}
\end{equation}
Similarly, for $1\le k\le p$, we have
\begin{equation*}
\begin{aligned}
     a_nb_k -b_na_k = c_1c_2(\gamma_1-\gamma_2)\begin{bmatrix}-\gamma_2^n &\gamma_1^n\end{bmatrix} \begin{bmatrix}\gamma_1^k\\[2pt]\gamma_2^k\end{bmatrix}.
\end{aligned}  
\end{equation*}
Consequently,
\begin{equation}\label{aB-bA}
\begin{aligned}
     (a_nB^T-b_nA^T)\Sigma_Z^{-1}(a_nB-b_nA)&=c_1^2c_2^2(\gamma_1-\gamma_2)^2\begin{bmatrix}-\gamma_2^n &\gamma_1^n\end{bmatrix}T\begin{bmatrix}-\gamma_2^n \\[2pt]\gamma_1^n\end{bmatrix}\\
     &=c_1^2c_2^2(\gamma_1-\gamma_2)^2\left(T_{11}\gamma_2^{2n}+T_{22}\gamma_1^{2n}-(T_{12}+T_{21})\gamma_1^n\gamma_2^n\right),
\end{aligned}  
\end{equation}
where $T$ is the $2 \times 2$ matrix defined by
\[
T = \sum_{i=1}^p \sum_{j=1}^p
\begin{bmatrix} \gamma_1^i \\[2pt] \gamma_2^i \end{bmatrix}
(\Sigma_Z^{-1})_{ij}
\begin{bmatrix} \gamma_1^j & \gamma_2^j \end{bmatrix}.
\]
%where $T = \sum_{i,j=1}^p\begin{bmatrix}\gamma_1^i\\[2pt]\gamma_2^i\end{bmatrix}(\Sigma^{-1}_Z)_{ij}\begin{bmatrix}\gamma_1^j&\gamma_2^j\end{bmatrix}$.

Taking the limit $n \to \infty$ in \eqref{MMSE-Vn2}, we obtain
\begin{equation}\label{power-eqn-sk2-proof}
   \begin{aligned}
      \lim_{n\to\infty}\E[|X_n|^2]%&=\lim_{n\to\infty}\E[|V_n-\E[V_n\mid\widetilde{Y}_1^{n-1}]|^2] \\
      &=\lim_{n\to\infty}\frac{a_n^2+b_n^2+\|a_nD_{2,n-1}-b_nD_{1,n-1}\|^2 +(a_nB^T-b_nA^T)\Sigma_Z^{-1}(a_nB-b_nA)}{\det(\Delta_{2,n-1})}\\
      &\overset{(f)}{=}\frac{L_{11}\gamma_1^2 +L_{22}\gamma_2^2+2L_{12}\gamma_1\gamma_2}{\Delta \gamma_1^2\gamma_2^2}\\
      &=\frac{1}{\Delta}\left(\frac{L_{11}}{\gamma_2^2}+\frac{L_{22}}{\gamma_1^2}+\frac{2L_{12}}{\gamma_1\gamma_2}\right)\\
      &= P,
   \end{aligned}
\end{equation}
where $(f)$ follows from (\ref{det-sk2}), (\ref{aD-bD}) and \eqref{aB-bA}. 

Therefore, under the condition $b\neq 0$,  it follows from \eqref{mutual-info-sk2-proof} and \eqref{power-eqn-sk2-proof} that the desired results in \eqref{I-SK2-thm} and \eqref{power-eqn-sk2-thm} hold, completing the proof for Case A.
%This establishes the power constraint (\ref{power-eqn-sk2}) for the SK(2) coding scheme. Together with (\ref{mutual-info-sk2-proof}), this completes the proof for Case A.




We now establish the result for Case B. Recall that in this case we have $\gamma_1=\gamma_2=\gamma\in \mathbb{R}$ with $|\gamma|>1$, which yields $a_n=(2-n)\gamma^{n-1}$ and $b_n=(n-1)\gamma^{n-2}$ as in (\ref{Vn-case-b}). Following the same construction as in \eqref{Y-tlta-re}, we define \(\{\widetilde{Y}\}\) accordingly, which leads to
\begin{equation*}\label{d12-case-c}
\left\{
\begin{aligned}
d_{1,n} &=\left[(2-n)L(\gamma^{-1})+L^\prime(\gamma^{-1})\gamma^{-1} \right]  \gamma^{n-1}\\[2pt]
d_{2,n} &= \left[(n-1)L(\gamma^{-1}) -L^\prime(\gamma^{-1})\gamma^{-1}   \right]\gamma^{n-2},
\end{aligned}
\right.
\end{equation*}
where $L^\prime(\gamma^{-1})$ denotes the derivative of $L(z)$ evaluated at $z = \gamma^{-1}$. 

Consequently, the scalar products $D^T_{1,n}D_{1,n}$, $D^T_{2,n}D_{2,n}$ and $D^T_{1,n}D_{2,n}$ are computed as follows:
\begin{equation}\label{D12n-case-c}
\begin{aligned}
D^T_{1,n}D_{1,n}&= L^2(\gamma^{-1})\left(\gamma^2S^{(2)}_{n-1}+1\right) -2L(\gamma^{-1})L^\prime(\gamma^{-1})\left(\gamma S^{(1)}_{n-1} -\gamma^{-1}\right)-R_1\\
&~~~~+\frac{(L^\prime(\gamma^{-1}))^2\gamma^{-2}}{\gamma^2-1}(\gamma^{2n}-1),\\
D^T_{2,n}D_{2,n}&= L^2(\gamma^{-1})\gamma^{-2}S_n^{(2)} -2L(\gamma^{-1})L^\prime(\gamma^{-1})\gamma^{-3}S_n^{(1)} +\frac{(L^\prime(\gamma^{-1}))^2\gamma^{-4}}{\gamma^2-1}(\gamma^{2n}-1)-R_2,\\
D^T_{1,n}D_{2,n}&= L(\gamma^{-1})\left(L^\prime(\gamma^{-1})+L^\prime(\gamma^{-1})\gamma^{-2}-L(\gamma^{-1})\gamma\right)S_{n-1}^{(1)}-R_{12}\\
&~~~~+L(\gamma^{-1})L^\prime(\gamma^{-1})\gamma^{-4}\gamma^{2n}(n-1) - L^2(\gamma^{-1})\gamma S^{(2)}_{n-1} \\
&~~~~-\frac{(L^\prime(\gamma^{-1}))^2\gamma^{-3}}{\gamma^2-1}(\gamma^{2n}-1) - L(\gamma^{-1})L^\prime(\gamma^{-1})\gamma^{-2},
\end{aligned}
\end{equation}
where
\begin{equation}\label{Sn1-case-c}
\begin{aligned}
S^{(1)}_n \triangleq \sum_{k=1}^{n-1}k\gamma^{2k}=\frac{1}{\gamma^2-1}n\gamma^{2n} - \frac{\gamma^4}{(\gamma^2-1)^2}\gamma^{2(n-1)} + \frac{\gamma^2}{(\gamma^2-1)^2}
\end{aligned}
\end{equation}
and
\begin{equation}\label{Sn2-case-c}
   \begin{aligned}
      S^{(2)}_n &\triangleq \sum_{k=1}^{n-1}k^2\gamma^{2k}=\frac{1}{\gamma^2-1}n^2\gamma^{2n} - \frac{2\gamma^2}{(\gamma^2-1)^2}n\gamma^{2n} + \frac{\gamma^4(\gamma^2+1)}{(\gamma^2-1)^3}\gamma^{2(n-1)}-\frac{(\gamma^2+1)\gamma^2}{(\gamma^2-1)^3}.
   \end{aligned}
\end{equation}
Proceeding analogously to the derivation of (\ref{det-sk2-1}) and (\ref{det-sk2}), it follows from \eqref{D12n-case-c}, \eqref{Sn1-case-c} and \eqref{Sn2-case-c} that 
\begin{equation}\label{det-sk2-case-c}
   \begin{aligned}
   \det(\Delta_{2,n})&=\frac{L^4(\gamma^{-1})}{(\gamma^2-1)^4}\gamma^{4(n+1)} + O(|\gamma|^{2n}).
   \end{aligned}
\end{equation}
Applying this to \eqref{pf-sk2-mutual-info} and subsequently taking the limit as $n\to\infty$, leads to
\begin{equation}\label{pf-sk2-mutual-info-casec}
\begin{aligned}
\bar{I}(V;Y)&=\lim_{n\to\infty}\frac{1}{n}I(V_1^n;Y_1^n)\\
&= 2\log |\gamma|.
\end{aligned}
\end{equation}
% This equality provides the desired lower bound on the mutual information rate achieved by the SK(2) coding scheme in Case B.

We next evaluate the limit of the power usage as in \eqref{MMSE-Vn2}. To this end, observe that
$$
a_n d_{2,k} - b_nd_{1,k} = \gamma^{n-3}(H^{(2)}_n)^T\begin{bmatrix}
   k\\1
\end{bmatrix}
\gamma^k,
$$
where 
$$
H^{(2)}_n \triangleq \begin{bmatrix}L(\gamma^{-1}) & -\left(nL(\gamma^{-1}) +\gamma^{-1}L^\prime(\gamma^{-1})\right)\end{bmatrix}^T.
$$
This implies that 
\begin{equation}\label{aD2-bD1-case-c}
\hspace{-0.3cm}
   \begin{aligned}
      \|a_nD_{2,n-1} - b_nD_{1,n-1}  \|^2 &= \sum_{k=p+1}^{n-1}(a_nd_{2,k} - b_nd_{1,k})^2\\
      &=\gamma^{2(n-3)}\sum_{k=p+1}^{n-1}(H^{(2)}_n)^T\begin{bmatrix}
         k\\ 1
      \end{bmatrix}
      \begin{bmatrix}
         k&1
      \end{bmatrix}
      H^{(2)}_n \gamma^{2k}\\[2pt]
      &= (H^{(2)}_n)^T\begin{bmatrix}
         \sum_{k=p+1}^{n-1} k^2\gamma^{2k} & \sum_{k=p+1}^{n-1}k\gamma^{2k}\\[2pt]
         \sum_{k=p+1}^{n-1}k\gamma^{2k}& \sum_{k=p+1}^{n-1}\gamma^{2k}
      \end{bmatrix}
      H^{(2)}_n \gamma^{2(n-3)}\\[2pt]
      &= \gamma^{2(n-3)}(H^{(2)}_n)^T\begin{bmatrix}
         S_n^{(2)}-\sum_{k=1}^pk^2\gamma^{2k} & S_n^{(1)}-\sum_{k=1}^pk\gamma^{2k}\\[2pt]
         S_n^{(1)})-\sum_{k=1}^pk\gamma^{2k}& \frac{\gamma^{2n}-\gamma^{2(p+1)}}{\gamma^2-1}
      \end{bmatrix}
      H^{(2)}_n.
   \end{aligned}
\end{equation}

Then, taking the limit as $n\to\infty$ in \eqref{MMSE-Vn2} gives
\begin{equation}\label{power-eqn-case-c}
   \hspace{-0.5cm}
   \begin{aligned}
      \lim_{n\to\infty}\E[|X_n|^2] %&= \lim_{n\to\infty}\E[|V_n - \E[V_n|\widetilde{Y}_1^{n-1}] |^2]\\
      &= \lim_{n\to\infty} \frac{a_n^2 +b_n^2 +\|a_nD_{2,n-1} - b_nD_{1,n-1}  \|^2 +(a_nB^T-b_nA^T)\Sigma_Z^{-1}(a_nB-b_nA)}{\det(\Delta_{2,n-1})}\\
      &\overset{(g)}{=} \frac{(\gamma^2-1)^3(L^\prime(\gamma^{-1}))^2}{\gamma^4 L^4(\gamma^{-1})} + \frac{2(\gamma^2-1)^2}{\gamma}\frac{L^\prime(\gamma^{-1})}{L^3(\gamma^{-1})} +\frac{\gamma^4-1}{L^2(\gamma^{-1})},
   \end{aligned}
\end{equation}
where $(g)$ follows from the definitions of $a_n$ and $b_n$ together with (\ref{det-sk2-case-c}) and (\ref{aD2-bD1-case-c}).



It remains to show that the left-hand side of \eqref{power-eqn-sk2-thm}, in the limit where $\gamma_1\to\gamma$ and $\gamma_2=\gamma\in\mathbb{R}$, coincides with the right-hand side of \eqref{power-eqn-case-c}. More precisely,
\begin{equation}\label{power-eqn-equivalence-case-c}
\hspace*{-0.3cm}
\begin{aligned}
\lim_{\substack{\gamma_1\to\gamma \\ \gamma_2=\gamma\in\mathbb{R}}} \frac{1}{\Delta}\left(\frac{L_{11}}{\gamma_2^2} + \frac{L_{22}}{\gamma_1^2} + \frac{2L_{12}}{\gamma_1\gamma_2}\right) &= \frac{(\gamma^2-1)^3(L^\prime(\gamma^{-1}))^2}{\gamma^4 L^4(\gamma^{-1})} + \frac{2(\gamma^2-1)^2}{\gamma}\frac{L^\prime(\gamma^{-1})}{L^3(\gamma^{-1})}+\frac{\gamma^4-1}{L^2(\gamma^{-1})}.
\end{aligned}
\end{equation}
Together with \eqref{pf-sk2-mutual-info-casec} and \eqref{power-eqn-case-c}, this immediately establishes the desired results as in \eqref{I-SK2-thm} and \eqref{power-eqn-sk2-thm} for Case B. The proof of the identity (\ref{power-eqn-equivalence-case-c}) is deferred to Appendix \ref{sec:proof_of_ref_power_eqn_equivalence_case_c}.
\end{proof}


Combining Theorem~\ref{ISK1-thm} (the case $b=0$) and Theorem~\ref{main-theorem-sk2} (the case $b\neq0$), we characterize SK(2) capacity $C_\mathrm{SK2}(P)$, which in turn provides a lower bound on $C_\mathrm{FB}(P)$.
\begin{co}\label{main-col-sk2:label}
Suppose the noise $\{ Z_i \}$ is an $AR(p)$ Gaussian process with parameters $\beta_k$, $\vert \beta_k\vert<1$ for all $k=1,2,...,p$, namely, it has the power spectral density
\begin{equation*}
S_Z(e^{i\theta}) = \vert H_Z(e^{i\theta}) \vert^2 = \frac{1}{\vert L_Z(e^{i\theta}) \vert^2} = \frac{1}{\vert\prod_{k=1}^{p}(1+\beta_ke^{i\theta})\vert^2}.
\end{equation*}
Then, we have
\begin{equation*}
C_\mathrm{SK2}(P)=  C_\mathrm{SK1}(P)\vee C^\mathrm{g}_\mathrm{SK2}(P),
\end{equation*}
where $C_\mathrm{SK1}(P)$ and $C^\mathrm{g}_\mathrm{SK2}(P)$ are given by \eqref{I-SK1-thm} and \eqref{I-SK2-thm}, respectively.
\end{co}
\begin{proof}
The result follows immediately from the definition of $C_\mathrm{SK2}(P)$, Theorem~\ref{ISK1-thm} (the case $b=0$), and Theorem~\ref{main-theorem-sk2} (the case $b\neq0$).
\end{proof}
As the simplest application of Theorem \ref{main-theorem-sk2} and Corollary~\ref{main-col-sk2:label}, we consider the case of a flat noise spectrum, defined by 
\begin{equation}\label{AWGN-SDF}
   S_Z(e^{i\theta})\equiv 1.
\end{equation}
This spectral density corresponds to a stationary noise process of the form $Z_i = W_i,~i\ge 1$, where $\{W_i\}_{i=1}^{\infty}$ is a zero-mean, unit-variance white Gaussian process. This corresponds precisely to the AWGN channel (\ref{AWGN-channel}). 

For the case $b = 0$, it follows directly from Theorem~\ref{ISK1-thm} and Corollary~\ref{main-col-sk2:label} that the SK(2) coding scheme reduces to the SK(1) coding scheme and is therefore optimal for this channel, as noted in Remark~\ref{remark1:label}. For the case $b\neq 0$, the following corollary demonstrates that, for this channel, the optimal pair $(\gamma_1,\gamma_2)$ in the optimization problem (\ref{I-SK2-thm})-(\ref{power-eqn-sk2-thm}) satisfies $|\gamma_1\gamma_2|=\sqrt{P+1}$.  Consequently, the SK(2) coding scheme yields a broader family of optimal coding schemes for the AWGN channel \eqref{AWGN-channel}.

\begin{co}\label{co-awgn-sk2:label}
Consider the noise process $\{Z_i\}_{i=1}^\infty$ with the power spectral density (\ref{AWGN-SDF}). Then the feedback capacity of the AWGN channel $Y_i = X_i +Z_i,~i=1,2,\dots$, subject to the power constraint $P$, can be achieved via the SK(2) coding scheme defined in \eqref{SK2-coding-input} and \eqref{SK2-coding-message} for $b\neq 0$ with any admissible parameter pair $(\gamma_1,\gamma_2)$ satisfying all root constraints and $|\gamma_1\gamma_2|=\sqrt{P+1}$.
\end{co}

\begin{proof}
For the AWGN channel \eqref{AWGN-channel}, we have $L_Z(e^{i\theta})\equiv 1$. Applying Theorem~\ref{main-theorem-sk2}, the optimization problem formulated in (\ref{I-SK2-thm}), (\ref{power-eqn-sk2-thm}) and (\ref{main-thm-notations}) can be simplified. More precisely, the power constraint (\ref{power-eqn-sk2-thm}) reduces to
\begin{equation}\label{awgn-sk2-power-eqn}
\gamma^2_1 \gamma^2_2-1\le P,
\end{equation}
where $(\gamma_1,\gamma_2)$ denote the characteristic roots of the equation (\ref{sk2-char-eqn}) satisfying $|\gamma_1|\wedge|\gamma_2|>1$. The objective is to maximize 
$$
C^\mathrm{g}_\mathrm{SK2}(\gamma_1,\gamma_2)\triangleq\log  |\gamma_1\gamma_2|
$$
subject to (\ref{awgn-sk2-power-eqn}).
% It is clearly to see that 
% \begin{equation*}
% \bar{I}_\mathrm{SK2}(\gamma_1,\gamma_2)\le \frac{1}{2}\log(P+1)
% \end{equation*}
% with equality if and only if $|\gamma_1\gamma_2| = \sqrt{P+1}$, which corresponds to any pair $(\gamma_1,\gamma_2)$ such that  $|\gamma_1\gamma_2| = \sqrt{P+1}$. 

We now consider two cases.

\textbf{Case A:} $\gamma_1 = \overline{\gamma_2}$ and $|\gamma_2| = \gamma > 1$. Then $C^\mathrm{g}_\mathrm{SK2}(\gamma_1,\gamma_2)=2\log \gamma$, and the power constraint (\ref{awgn-sk2-power-eqn}) becomes $\gamma^4-1\le P$. Hence,
\begin{equation*}
C^\mathrm{g}_\mathrm{SK2}(\gamma_1,\gamma_2)\le \frac{1}{2}\log(P+1)
\end{equation*}
with equality if and only if $\gamma = (P+1)^{1/4}$, i.e., $\gamma_1=\overline{\gamma_2}$ and $|\gamma_2|= (P+1)^{1/4}$. 

\textbf{Case B:} $\gamma_1, \gamma_2 \in \mathbb{R}$ with $\gamma_1 \neq \gamma_2$. From (\ref{awgn-sk2-power-eqn}), we obtain
\begin{equation*}
\begin{aligned}
\sup_{\substack{\gamma_1,\gamma_2\in\mathbb{R}\\\gamma_1\neq\gamma_2}} C^\mathrm{g}_{\mathrm{SK2}}(\gamma_1,\gamma_2) &= \sup \log |\gamma_1\gamma_2|= \frac{1}{2}\log(P+1),
\end{aligned}
\end{equation*}
where the supremum is attained by any pair $(\gamma_1,\gamma_2)$ satisfying $|\gamma_1\gamma_2|=\sqrt{P+1}$.

Combining both cases, we have shown that the maximum $C^\mathrm{g}_\mathrm{SK2}(P)$ in \eqref{I-SK2-thm} is achieved whenever $(\gamma_1,\gamma_2)$ satisfy $|\gamma_1\gamma_2|=\sqrt{P+1}$. In particular, the well-known feedback capacity of the AWGN channel, $\frac{1}{2}\log(P+1)$, is achieved by the SK(2) coding scheme for $b\neq 0$ with any such pair $(\gamma_1,\gamma_2)$ satisfying $|\gamma_1\gamma_2|=\sqrt{P+1}$. This completes the proof.
\end{proof}

We next consider the stationary AR(1) Gaussian channel 
\[
Y_i=X_i+Z_i,\quad i\ge1,
\]
with noise power spectral density
$$
S_Z(e^{i\theta})=\frac{1}{|1+\beta_1e^{i\theta}|^2},\qquad \beta_1 \in (-1,1).
$$
This spectral density corresponds to the stationary noise process
$$
Z_i +\beta_1 Z_{i-1} = W_i,\quad i\in\mathbb{Z}.
$$ 
By \cite[Theorem 5.2]{su2024feedback}, the feedback capacity of this channel under power constraint $P$ is
%According to \cite{kim2009feedback} and \cite{yang2005feedback}, 
\begin{equation}\label{CFB-AR1}
C_\mathrm{FB}(P) = -\log x_0,
\end{equation}
where $x_0$ is the unique positive root of the fourth-order polynomial equation
\begin{equation}\label{Power-Eqn-AR1}
Px^2 = \frac{1-x^2}{(1+ |\beta_1| x)^2}.
\end{equation}
As noted in Remark~\ref{remark1:label}, the SK(1) coding scheme defined in (\ref{SK1-coding-X}) and (\ref{SK1-coding-message}) is optimal for this channel; specifically,
\[
   C_\mathrm{SK1}(P)=C_\mathrm{FB}(P).
\]
Hence, by Corollary~\ref{main-col-sk2:label},
\begin{equation}\label{CSK2-AR1_equal}
   C_\mathrm{SK2}(P)=C_\mathrm{FB}(P).
\end{equation}
Thus, the SK(2) coding scheme with $b=0$ is optimal for the stationary AR(1) Gaussian channel. The following example shows, however, that for the case $b\neq0$ the SK(2) coding scheme may not be optimal unless $\beta_1=0$  (i.e., the AWGN channel \eqref{AWGN-channel}).
\begin{example}\label{example-ar1:label}
%Equivalently, it is attained by the SK(1) scheme introduced in this paper, satisfying
%\begin{equation}\label{ISK1-AR1-CFB}
%I_{SK1}(P)=C_{FB}(P).
%\end{equation}
From Theorem~\ref{main-theorem-sk2} together with (\ref{CFB-AR1}) and \eqref{Power-Eqn-AR1}, we obtain
$$
C^\mathrm{g}_{\mathrm{SK2}}(P) < C_{\mathrm{FB}}(P),
$$
for all $\beta_1\in(-1,1)\setminus \{0\}$. For different values of $\beta_1$ and $P = 1, 5$, Figure~\ref{fig:ISK-AR1-3P-diff} plots $C_{\mathrm{SK1}}(P)$, $C^\mathrm{g}_{\mathrm{SK2}}(P)$, and their difference as functions of $\beta_1$.
\begin{figure*}[htbp!]
    \centering
    %\hspace*{-1.5cm}
    \centerline{\includegraphics[width=17cm]{AR1_P15.jpg}}
    \caption{Comparison of $C_\mathrm{SK1}(P)$ and $C^\mathrm{g}_\mathrm{SK2}(P)$ for stationary AR(1) Gaussian channels with $P=1,5$.}
    \label{fig:ISK-AR1-3P-diff}
\end{figure*}
\end{example}
Therefore, we have established the optimality of the SK(2) coding scheme for stationary AR(1) Gaussian channels, as summarized in the following corollary.
\begin{co}\label{co-ar1-sk2:label}
The SK(2) coding scheme achieves the feedback capacity of stationary AR(1) Gaussian channels.
\end{co}
\begin{proof}
The result follows directly from Corollary~\ref{main-col-sk2:label}, Example~\ref{example-ar1:label} and \eqref{CSK2-AR1_equal}.
\end{proof}

We next examine the performance of the SK(2) coding scheme over stationary AR(2) Gaussian channels. We first present numerical comparison for a general AR(2) spectrum and then give an analytical family for which the genuine SK(2) coding scheme achieves feedback capacity and strictly outperforms the SK(1) coding scheme.

%We now present a second example, which examines stationary AR(2) Gaussian channels and demonstrates that the SK(2) coding scheme can achieve strictly higher rates than the SK(1) coding scheme in certain scenarios.
\begin{example}\label{example-ar2:label}
Consider a stationary AR(2) Gaussian channel with noise power spectral density
\begin{equation*}\label{AR2-SDF}
S_Z(e^{i\theta})=\frac{1}{|(1+\beta_1e^{i\theta})(1+\beta_2e^{i\theta})|^2},\qquad |\beta_1|,|\beta_2|<1.
\end{equation*}
This spectrum corresponds to the stationary noise process 
$$
Z_i +(\beta_1+\beta_2) Z_{i-1} +\beta_1\beta_2 Z_{i-2} = W_i,\quad i\in\mathbb{Z},
$$
where $\{W_i\}$ is a white Gaussian process with zero mean and unit variance. For the AR coefficients to be real, the pair $(\beta_1,\beta_2)$ must satisfy $\beta_1+\beta_2\in \mathbb{R}$ and $\beta_1\beta_2\in\mathbb{R}$ (i.e., $\beta_1$ and $\beta_2$ are either real or complex conjugates). 

Under the power constraint $P$, Theorem~\ref{ISK1-thm} yields
\begin{equation*}\label{ISK1-AR2}
C_{\mathrm{SK1}}(P)= \max\log |x_0|,
\end{equation*}
where the maximum is taken over all real roots of the polynomial equation
\begin{equation*}\label{Power-Eqn-AR2}
P = \frac{x^2-1}{\left((1+ \beta_1 x^{-1})(1 + \beta_2 x^{-1})\right)^2}.
\end{equation*}

%Numerically evaluating the optimization problems in Theorem~\ref{main-theorem-sk2} reveals nonempty parameters ranges for which

Theorem~\ref{main-theorem-sk2}, together with some straightforward numerical computations (see Figure~\ref{fig:ISK-AR2-3Beta}), shows that in certain cases
\[
C^\mathrm{g}_{\mathrm{SK2}}(P) > C_{\mathrm{SK1}}(P).
\]
By Corollary~\ref{main-col-sk2:label}, this implies
\begin{equation}\label{CISK2-ar2}
C_\mathrm{SK2}(P)= C^\mathrm{g}_\mathrm{SK2}(P)>C_{\mathrm{SK1}}(P)=R_\mathrm{Butman}(P).
\end{equation}
\begin{figure}[htbp!]
    \centering
    %\hspace*{-1.5cm}
    \centerline{\includegraphics[width=17cm]{AR2_beta59_P1.jpg}}
    \caption{Comparison of $C_\mathrm{SK1}(P)$ and $C^\mathrm{g}_\mathrm{SK2}(P)$ for stationary AR(2) Gaussian channels with $\beta_2=0.5,~0.9$ and $P=1$.}
    \label{fig:ISK-AR2-3Beta}
\end{figure}
\end{example}
Thus, the numerical results illustrate that the advantage of second-order message recursion occurs for a nontrivial range of AR(2) Gaussian channels. The analytical separation needed to disprove Conjecture~\ref{Butman-conj:label} is established below. A particularly transparent analytical family is obtained by removing the first-order noise term. In this case, the even- and odd-indexed noise samples form two independent, interleaved AR(1) processes. This structure is naturally matched by a second-order message recursion. The following proposition shows that SK(2) not only outperforms SK(1) for this family, but also achieves feedback capacity.



%The following example illustrates a class of AR(2) Gaussian channels for which the SK(2) coding scheme attains the feedback capacity, whereas the SK(1) and Butman's coding scheme are strictly sub-optimal.

% \begin{example}\label{Interleaved-ar2-capacity:label}
% Consider the AR(2) Gaussian noise process
% \begin{equation*}
% Z_i + \beta Z_{i-2} = W_i, \qquad i \in \mathbb{Z},
% \end{equation*}
% which decomposes into two independent and interleaved AR(1) processes, one on the even indices and one on the odd indices. This channel therefore admits the same feedback capacity as the AR(1) channel parameter $\beta$, and the feedback capacity can be attained by signaling separately for each interleaved AR(1) channel.

% Under the average power constraint $P$, the feedback capacity is given by
% \begin{equation}
% C_{\mathrm{FB}}(P) = \log x_0,
% \end{equation}
% where $x_0$ is the unique positive root of the fourth-order polynomial equation
% \begin{equation}
% P = \frac{x^2 - 1}{\bigl(1 + |\beta| x^{-1}\bigr)^2}.
% \end{equation}

% Applying Theorems~\ref{ISK1-thm} and~\ref{main-theorem-sk2}, we obtain
% \[
% C_{\mathrm{SK2}}(P) = \bar{I}_{\mathrm{SK2}}(P) = C_{\mathrm{FB}}(P),
% \]
% with the optimal parameter pair \((\gamma_1, \gamma_2)\) in \eqref{I-SK2-thm} and \eqref{power-eqn-sk2-thm} given by
% \[
% (\gamma_1, \gamma_2) = 
% \begin{cases}
% (\sqrt{x_0}, -\sqrt{x_0}), & \beta > 0, \\[2pt]
% (i \sqrt{x_0}, -i \sqrt{x_0}), & \beta < 0,
% \end{cases}
% \]
% where $i = \sqrt{-1}$. Thus, for this class of AR(2) channels, the SK(2) coding scheme achieves the feedback capacity.
% \end{example}

\begin{pr}
\label{prop:interleaved-AR1}
Consider the stationary AR(2) Gaussian channel
$$
Y_i=X_i+Z_i,\qquad Z_i+\beta Z_{i-2}=W_i,\quad i\in\mathbb Z,
$$
where \(0<|\beta|<1\) and \(\{W_i\}\) is a zero-mean, unit-variance white Gaussian process. Then, for every $P>0$,
$$
C_{\mathrm{SK2}}(P) = C_{\mathrm{FB}}(P) > C_{\mathrm{SK1}}(P) = R_{\mathrm{Butman}}(P).
$$
Consequently, Conjecture~\ref{Butman-conj:label} is false for stationary AR(2) Gaussian channels.
\end{pr}

\begin{proof}
Define $Z_k^{(e)}=Z_{2k}$ and $Z^{(o)}_k=Z_{2k-1}$ for $k\in\mathbb{Z}$. Since the noise process $\{Z_i\}$ satisfies $Z_i+\beta Z_{i-2}=W_i$, we have
$$
Z^{(e)}_{k}+\beta Z_{k-1}^{(e)}=W_{2k}, \qquad Z^{(o)}_{k}+\beta Z^{(o)}_{k-1}=W_{2k-1}.
$$
Therefore, $\{Z^{(e)}_k\}$ and $\{Z^{(o)}_k\}$ are two independent stationary AR(1) Gaussian processes with parameter $\beta$, driven by the independent white Gaussian noises $\{W_{2k}\}$ and $\{W_{2k-1}\}$, respectively. Consequently, the original AR(2) Gaussian channel can be viewed as two independent parallel AR(1) Gaussian channels obtained by separating the even and odd channel uses.

Let $C_{\mathrm{FB}}(P;\beta)$ denote the feedback capacity of either AR(1) subchannel. Let $P_e$ and $P_o$ denote the average input powers per use of the even and odd AR(1) subchannels, respectively. Since each subchannel occupies one half of the original channel uses, the overall average-power constraint becomes
$$
\frac{P_e+P_o}{2}\le P.
$$
Consequently, the feedback capacity of the original channel is
\begin{equation*}\label{CFB-interleave-ar2}
\begin{aligned}
C_{\mathrm{FB}}(P)
&= \max_{\frac{P_{\mathrm e}+P_{\mathrm o}}{2}\le P}
\left\{
\frac12 C_{\mathrm{FB}}(P_{\mathrm e};\beta) + \frac12 C_{\mathrm{FB}}(P_{\mathrm o};\beta) \right\}.
\end{aligned}
\end{equation*}
Since $C_\mathrm{FB}(P;\beta)$ is concave and non-decreasing in $P$, we have
$$
\frac12 C_\mathrm{FB}(P_{\mathrm e};\beta) + \frac12 C_{\mathrm{FB}}(P_{\mathrm o};\beta)
\le
C_{\mathrm{FB}}\left(\frac{P_{\mathrm e}+P_{\mathrm o}}{2};\beta\right)
\le
C_{\mathrm{FB}}(P;\beta),
$$
with equality attained by $P_{\mathrm e}=P_{\mathrm o}=P$. Hence
\begin{equation}\label{CFB-interleave-ar2-2}
C_{\mathrm{FB}}(P)= C_{\mathrm{FB}}(P;\beta)=\log x_0,
\end{equation}
where $x_0>1$ is the unique solution of
\begin{equation}\label{eq:interleaved-capacity}
P= \frac{x_0^2-1} {\left(1+|\beta|x_0^{-1}\right)^2}.
\end{equation}

Now, applying the SK(2) coding scheme \eqref{SK2-coding-input}--\eqref{SK2-coding-message} and choosing the characteristic pair $(\gamma_1,\gamma_2)$ as
\begin{equation}\label{chara-pair-interleave-ar2}
(\gamma_1,\gamma_2) =
\begin{cases}
\left(\sqrt{x_0},-\sqrt{x_0}\right), & \beta>0,\\[1mm]
\left(\mathrm{i}\sqrt{x_0},-\mathrm{i}\sqrt{x_0}\right), & \beta<0.
\end{cases}
\end{equation}
Then
$$
\gamma_1+\gamma_2=0, \qquad \gamma_1\gamma_2=-\operatorname{sign}(\beta)x_0.
$$
In particular, 
$$
|\gamma_1|=|\gamma_2| = \sqrt{x_0}>1,
$$
and both $\gamma_1+\gamma_2$ and $\gamma_1\gamma_2$ are real. Moreover, 
$$
b = -\gamma_1\gamma_2=\operatorname{sign}(\beta)x_0\neq 0.
$$
Thus, it remains to verify that the characteristic pair $(\gamma_1,\gamma_2)$ satisfies the power constraint in Theorem~\ref{main-theorem-sk2}.
%and the corresponding message recursion is
% $$
% V_{i+1}=\operatorname{sign}(\beta)x_0V_{i-1}.
% $$
% Thus, defining $V_k^{(o)}=V_{2k-1}$ and $V_k^{(e)}=V_{2k}$ for $k\ge 1$, we have
% $$
% \begin{cases}
% V_{k+1}^{(o)} = \operatorname{sign}(\beta)x_0V_{k}^{(o)}\\[4pt]
% V_{1}^{(o)} = U_1
% \end{cases},
% \qquad
% \begin{cases}
% V_{k+1}^{(e)} = \operatorname{sign}(\beta)x_0V_{k}^{(e)}\\[4pt]
% V_{1}^{(e)} = U_2
% \end{cases}.
% $$
%We next verify that this the characteristic pair $(\gamma_1,\gamma_2)$ satisfies the power constraint in Theorem~\ref{main-theorem-sk2}.
For both $\beta>0$ and $\beta<0$, 
$$
\gamma_1^2=\gamma_2^2=\operatorname{sign}(\beta) x_0,
$$
and, since $L_Z(z)=1+\beta z^2$,
$$
L_Z(\gamma_1^{-1})=L_Z(\gamma_2^{-1}) = 1 +\frac{\beta}{\operatorname{sign}(\beta) x_0}=1 +\frac{|\beta|}{ x_0}.
$$
Moreover,
$$
c_1=c_2=\frac{1}{2\operatorname{sign}(\beta) x_0}.
$$
Substituting these expressions into the power constraint \eqref{power-eqn-sk2-thm} gives
$$
\frac{x_0^2-1}{(1 +|\beta| x^{-1}_0)^2}=P,
$$
where the equality follows from \eqref{eq:interleaved-capacity}. Thus, the chosen characteristic pair $(\gamma_1,\gamma_2)$ in \eqref{chara-pair-interleave-ar2} is feasible in Theorem~\ref{main-theorem-sk2}. Consequently,
$$
C^\mathrm{g}_\mathrm{SK2}(P)\ge \log|\gamma_1\gamma_2|=\log x_0.
$$
Since $C_\mathrm{SK2}(P)\ge C^\mathrm{g}_\mathrm{SK2}(P)$, together with \eqref{CFB-interleave-ar2-2}, it holds that
$$
 C_{\mathrm{FB}}(P)=\log x_0 \le C_{\mathrm{SK2}}(P) \le C_{\mathrm{FB}}(P),
$$
and hence
\begin{equation}\label{CSK2-CFB-AR2-Inter}
C_\mathrm{SK2}(P)=C_\mathrm{FB}(P) = \log x_0.
\end{equation}

It remains to prove that $C_{\mathrm{SK1}}(P)<C_{\mathrm{FB}}(P)$. To see this, Theorem~\ref{ISK1-thm} gives
\begin{equation}\label{CFB-sk1-inter}
C_\mathrm{SK1}(P)=\max\log |\gamma|,
\end{equation}
where the maximum is taken over all $\gamma\in\mathbb{R}$ with $|\gamma|>1$ satisfying 
\begin{equation}\label{eq:SK1-interleaved}
P= \frac{\gamma^2-1} {\left(1+\beta \gamma^{-2}\right)^2}.
\end{equation}
Let $\gamma_0$ be an optimizer of \eqref{CFB-sk1-inter} and define
$$
\Phi_\beta(x) = \frac{x^2-1}{\left(1+|\beta| x^{-1}\right)^2}, \qquad x>1.
$$
Since
$$
\Phi_\beta'(x) = \frac{2x\left(x^3+2|\beta| x^2-|\beta|\right)} {(x+|\beta|)^3} >0, \qquad x>1,
$$
the function $\Phi_\beta$ is strictly increasing. Moreover, 
$$
1+\frac{|\beta|}{|\gamma_0|} > 1+\frac{\beta}{\gamma_0^2} >0.
$$
Using \eqref{eq:SK1-interleaved}, we obtain
$$
\Phi_\beta(|\gamma_0|) < \frac{\gamma_0^2-1} {\left(1+\beta \gamma_0^{-2}\right)^2} = P.
$$
On the other hand, \eqref{eq:interleaved-capacity} gives
$$
\Phi_\beta(x_0)=P.
$$
The strict monotonicity of $\Phi_\beta$ implies $|\gamma_0|<x_0$. Consequently
$$
C_{\mathrm{SK1}}(P) = \log |\gamma_0| < \log x_0 = C_{\mathrm{FB}}(P),
$$
as desired.

Finally, it follows from $R_{\mathrm{Butman}}(P) = C_{\mathrm{SK1}}(P)$ and \eqref{CSK2-CFB-AR2-Inter} that 
$$
R_{\mathrm{Butman}}(P) < C_{\mathrm{SK2}}(P) = C_{\mathrm{FB}}(P),
$$
which disproves Conjecture~\ref{Butman-conj:label}. This completes the proof. 
\end{proof}


\section{Conclusion and Outlook}

We introduced the SK(2) class of Gaussian feedback coding schemes for stationary AR($p$) Gaussian channels, in which the message process evolves according to a second-order deterministic recursion. We derived a closed-form characterization of its maximal achievable rate and showed that the SK(1) subclass provides a branch-complete and sign-consistent formulation of Butman's equal-energy construction. Since SK(2) includes SK(1) as a special case, it recovers the known feedback capacities of the AWGN and AR(1) Gaussian channels. More importantly, for certain AR(2) Gaussian channels, a genuinely second-order recursion achieves a rate strictly larger than that of SK(1); for an interleaved AR(1) subclass, it achieves feedback capacity.

These results show that the first-order SK paradigm is not the end of the story for Gaussian channels with memory. Although first-order SK/Butman signaling is optimal for AWGN and AR(1) noise, higher-order message recursions can exploit more complex noise memory and unlock achievable rates that are inaccessible to first-order schemes. In particular, the order of the signaling recursion emerges as an essential design dimension rather than merely a technical extension of the classical SK construction.

The feedback capacity of general AR($p$) Gaussian channels, however, remains open. A natural direction is therefore to investigate SK($q$) schemes based on deterministic message recursions of order $q>2$ and to determine whether increasing the recursion order yields a convergent hierarchy of achievable rates. More broadly, it remains to characterize the relationship between the order and structure of the noise spectrum and the minimal message-recursion order required for capacity-achieving. Resolving these questions may provide a constructive route toward the feedback capacity of general finite-order ARMA Gaussian channels.


%\textbf{Acknowledgment}: 
\section*{Appendices}
\appendix

\section{Proof of Lemma \ref{lemma1}}\label{sec:proof_of_lemma_ref_lemma1}
For any nonzero $x\in \mathbb{R}^{n}$, since $S$ is symmetric positive definite, we have 
\begin{equation*}
\begin{aligned}
x^T (S + U U^T)x &= x^T S x + x^T U U^T x\\
&=x^T S x + \|U^T x  \|^2\\
&\geq x^T S x\\
&>0,
\end{aligned}
\end{equation*}
which shows that $S + U U^T$ is positive-definite and hence invertible. Then, we use the Woodbury matrix identity \cite{higham2002accuracy} to establish
\begin{equation*}
\begin{aligned}
(S +UU^T )^{-1} &= S^{-1} -S^{-1}U(I_m+U^T S^{-1} U)^{-1}U^T S^{-1}\\
&=S^{-1} - S^{-1}UA^{-1}U^T S^{-1}
\end{aligned}
\end{equation*}
and thus
\begin{equation*}
\begin{aligned}
U^T (S+U U^T)^{-1}U &= U^T S^{-1} U - U^T S^{-1} U A^{-1} U^T S^{-1}U\\
&=A-I_m -(A-I_m) A^{-1}(A-I_m)\\
&= I_m-A^{-1}.
\end{aligned}
\end{equation*}
This further immediately implies that for $1\le i,j\le m$
\begin{equation}\label{lemma-eqn-1}
u_i^T (S +UU^T)^{-1}u_j = \delta_{ij} - (A^{-1})_{ij},
\end{equation}
where $\delta_{ij}$ is the Kronecker delta function. Furthermore, since 
$$
A^{-1} = \frac{1}{\det(A)}\text{adj}(A),
$$
where \text{adj}$(A)$ is the adjugate matrix of $A$, it holds that 
$$
(A^{-1})_{ij} = (-1)^{i+j}\frac{1}{\det(A)} M_{ji},
$$
which, together with $(\ref{lemma-eqn-1})$, immediately completes this proof, as desired.


\section{Proof of Theorem \ref{ISK1-thm}}\label{sec:proof_of_theorem_ref_isk1_thm}
Consider the SK(1) coding scheme formulated by \eqref{SK1-coding-X} and \eqref{SK1-coding-message} over the stationary AR($p$) Gaussian channel \eqref{channel-arp} with the noise spectral density function $S_Z(e^{i\theta})$ given by \eqref{sdf-arp}. It follows directly from \eqref{SK1-coding-message} that
\[
   V_n=\gamma^{n-1}U.
\]
Following the same arguments as used in \eqref{Y-tlta}, \eqref{Y-tlta-2} and \eqref{d1d2}, we can express
\begin{equation}
   \widetilde{Y}_n = 
\begin{cases} 
\gamma^{n-1} U + Z_n, & \text{for } n \le p; \\[5pt]
d_n U + W_n, & \text{for } n \ge p+1,
\end{cases}
\end{equation}
where 
\begin{equation}\label{SK1-d}
   d_n = \gamma^{n-1}L(\gamma^{-1}).
\end{equation}
Thus, the mutual information between $V_1^n$ and $Y_1^n$ can be expressed as 
\begin{equation}\label{mutual-info-rate-sk1}
   \begin{aligned}
      \frac{1}{n}I(V_1^n;Y_1^n)&= \frac{1}{n}I(U;Y_1^n)\\
      &=\frac{1}{n}I(U;\widetilde{Y}_1^n)\\
      &=-\frac{1}{2n}\log\E[|U-\E[U\mid\widetilde{Y}_1^n]|^2].
   \end{aligned}
\end{equation}
Define $G =  \begin{bmatrix}
         1 & \gamma &\cdots&\gamma^{p-1}
      \end{bmatrix}^T$, $D_n = \begin{bmatrix}
         d_{p+1} & d_{p+2}&\cdots&d_n
      \end{bmatrix}^T$ and $\widetilde{D}_n = \begin{bmatrix} G^T &D_n^T
      \end{bmatrix}^T$.
Since $U$ and $\widetilde{Y}_1^n$ are jointly Gaussian, letting $\Sigma_Z=\mathrm{Var}(Z_1^p)$, we obtain
\begin{equation}\label{mmse-sk1-pf}
   \begin{aligned}
      \E[|U-\E[U\mid\widetilde{Y}_1^n|^2] &= 1 - \begin{bmatrix}G^T &D_n^T \end{bmatrix}\begin{bmatrix}\Sigma_{Z}+GG^T&GD_n^T\\[2pt]D_nG^T&I+D_nD_n^T \end{bmatrix}^{-1}\begin{bmatrix}G \\[2pt]D_n\end{bmatrix}\\
      &=1 - \widetilde{D}_n^T\left(\begin{bmatrix}\Sigma_Z &0\\[2pt]0 &I\end{bmatrix}+\widetilde{D}_n\widetilde{D}_n^T\right)^{-1}\widetilde{D}_n \\
      &\overset{(a)}{=}  \frac{1}{1+G^T\Sigma_Z^{-1}G+D_n^TD_n}\\
      &= \frac{1}{1+G^T\Sigma_Z^{-1}G - L^2(\gamma^{-1})G^TG +\frac{L^{2}(\gamma^{-1})}{\gamma^{2}-1}(\gamma^{2n}-1)},
   \end{aligned}
\end{equation}
where $(a)$ follows from Lemma \ref{lemma1}.  Substituting this result into \eqref{mutual-info-rate-sk1} and taking the limit as $n\to\infty$, we obtain
\begin{equation}\label{proof-sk1-mr}
\begin{aligned}
\lim_{n\to\infty}\frac{1}{n}I(V_1^n;Y_1^n) = \log|\gamma|.
\end{aligned}
\end{equation}
Moreover, following a similar approach as in \eqref{power-eqn-sk2}, we have
\begin{equation}\label{power-eqn-sk1-pf}
   \begin{aligned}
   \E[|X_n|^2] &=\E[|V_n - \E[V_n\mid Y_1^{n-1}]|^2]\\
            &= \E[|V_n - \E[V_n\mid \widetilde{Y}_1^{n-1}]|^2]\\
            &=\gamma^{2(n-1)}\E[|U- \E[U\mid\widetilde{Y}_{1}^{n-1}]|^2].
   \end{aligned}
\end{equation}
Substituting \eqref{mmse-sk1-pf} into \eqref{power-eqn-sk1-pf} and letting $n\to\infty$  yields
\begin{equation}\label{proof-sk1-power}
   \begin{aligned}
   \lim_{n\to\infty}\E[|X_n|^2] &= \frac{\gamma^2-1}{L^2(\gamma^{-1})}= P.
 \end{aligned}
\end{equation}
Consequently, to determine $C_\mathrm{SK1}(P)$, one must maximize the mutual information rate in \eqref{proof-sk1-mr} subject to the power constraint \eqref{proof-sk1-power}. This completes the proof, as desired.

\section{Counterexamples for Butman's AR(p) Derivation}\label{sec:example_of_ar_2_gaussian_channel}
The following two examples illustrate the sign-compatibility and branch issues discussed in Section~\ref{subsec:relation_to_butman}.
\begin{example}[Failure of sign compatibility]\label{Butman-example-1:label}
Let $0<t<1$ and take
\[
    \beta_1=t,\qquad \beta_2=-t.
\]
Then
\[
    L_Z(z)=(1+tz)(1-tz)=1-t^2z^2.
\]
Hence
$$
    B(x)=-t^2x^{-2}<0,\qquad x\neq0 .
$$

For any \(P>0\), \eqref{butman-abs-eqn} becomes
\begin{equation}\label{Butman-Counterexample-pe-1}
    x^2 = 1+P\left(1+t^2x^{-2}\right)^2.
\end{equation}
Clearly, every real root of \eqref{Butman-Counterexample-pe-1} satisfies $|x|>1$. Let $r=|x|$. Under the equal-energy condition $x_i=x$, it follows from \eqref{Butman-ee-cond} that
\[
    s_i=\operatorname{sign}(x)^2s_{i-2}=s_{i-2}.
\]
On the other hand, since
\[
    S_i(r)=b_2s_{i-2}r^{-2}=-t^2s_{i-2}r^{-2},
\]
it holds that
\[
    \operatorname{sign}(S_i(r))=-s_{i-2}.
\]
Thus, the sign-selection rule \eqref{Butman-sign-cond} requires
\[
    s_i=-s_{i-2},
\]
which contradicts the equal-ratio requirement \(s_i=s_{i-2}\). This contradiction is independent of the sign of $x$. Therefore, no real root of \eqref{butman-abs-eqn} can be implemented together with both the sign-selection rule and the equal-ratio condition.

For completeness, the equation \eqref{Butman-Counterexample-pe-1} indeed has real roots. Setting \(u=x^2\), the equation \eqref{Butman-Counterexample-pe-1} becomes
\[
    u=1+P\left(1+\frac{t^2}{u}\right)^2 .
\]
Let
\[
    f(u)=u-1-P\left(1+\frac{t^2}{u}\right)^2.
\]

Then, \(f(1)<0\), \(f(u)\to\infty\) as \(u\to\infty\), and
\[
    f^\prime(u) = 1+\frac{2Pt^2}{u^2}\left(1+\frac{t^2}{u}\right)>0 .
\]
Consequently, there is a unique \(u_\ast>1\), and the real roots are \(x=\pm\sqrt{u_\ast}\). Both are incompatible with the recursive sign selection, as shown above.
\end{example}

\begin{example}[Failure of fixed-sign absolute-value removal]
\label{Butman-example-2:label}
Let \(0<t<1\) and take $ \beta_1=ti,~\beta_2=-ti$.
Then
$$
    L_Z(z)=(1+tiz)(1-tiz)=1+t^2z^2.
$$
Consequently,
$$
    B(x)=t^2x^{-2}>0,\qquad x\in\mathbb R,\ x\neq0 .
$$
Hence
\[
    |B(x)|=B(x)\neq -B(x)
\]
for every real \(x\neq0\). Therefore, the passage from the absolute-value expression \eqref{butman-abs-eqn} to Butman's expression \eqref{butman-final-eqn-1} is not justified for this AR(2) Gaussian channel.
\end{example}


\section{Proof of (\ref{power-eqn-equivalence-case-c})}\label{sec:proof_of_ref_power_eqn_equivalence_case_c}
For notational simplicity, let 
\[
   f(\gamma_1,\gamma_2)= \frac{1}{\Delta}\left(\frac{L_{11}}{\gamma_2^2} + \frac{L_{22}}{\gamma_1^2} + \frac{2L_{12}}{\gamma_1\gamma_2}\right).
\]
It then follows that 
\begin{equation}\label{f}
   f(\gamma_1,\gamma_2) = \frac{(\gamma_1^2-1)(\gamma_2^2-1)(\gamma_1\gamma_2-1)^2}{(\gamma_1 - \gamma_2)^2L^2(\gamma_1^{-1})L^2(\gamma_2^{-1})}g(\gamma_1,\gamma_2),
\end{equation}
where 
\[
   g(\gamma_1,\gamma_2)=\frac{L^2(\gamma_1^{-1})}{\gamma_1^2-1} + \frac{L^2(\gamma_2^{-1})}{\gamma_2^2-1}-\frac{2L(\gamma_1^{-1})L(\gamma_2^{-1})}{\gamma_1\gamma_2-1}.
\]
Let $\gamma_2 = \gamma \in \mathbb{R}$. The Taylor expansion of $g(\gamma_1, \gamma)$ in $\gamma_1$ around $\gamma_1 = \gamma$ is:
%Let $\gamma_2=\gamma\in\mathbb{R}$ and by Taylor's expansion of $g(\gamma_1,\gamma)$ as a function of $\gamma_1\in\mathbb{R}$:
\begin{equation}\label{taylor-g}
   \begin{aligned}
      g(\gamma_1,\gamma)= \sum_{n=0}^\infty\frac{1}{n!}\frac{\partial^ng}{\partial\gamma_1^n}(\gamma,\gamma)(\gamma_1-\gamma)^n
   \end{aligned}
\end{equation}
Clearly,
\begin{equation}\label{g-0}
   g(\gamma,\gamma)=0
\end{equation}
and
\begin{equation*}
   \begin{aligned}
      \frac{\partial g}{\partial \gamma_1}(\gamma_1,\gamma)= -\frac{2L(\gamma_1^{-1})L^\prime(\gamma_1^{-1})}{\gamma_1^2(\gamma_1^2-1)} -\frac{2\gamma_1L^2(\gamma_1^{-1})}{(\gamma_1^2-1)^2} + \frac{2L(\gamma^{-1})L^\prime(\gamma_1^{-1})}{\gamma_1^2(\gamma_1\gamma-1)}+\frac{2\gamma L(\gamma^{-1})L(\gamma_1^{-1})}{(\gamma_1\gamma-1)^2},
   \end{aligned}
\end{equation*}
which implies that 
\begin{equation}\label{g-1}
   \frac{\partial g}{\partial \gamma_1}(\gamma,\gamma)=0.
\end{equation}
Furthermore, we have 
\begin{equation*}
   \begin{aligned}
      \frac{\partial^2 g}{\partial \gamma^2_1}(\gamma_1,\gamma)&= \frac{2(L^\prime(\gamma_1^{-1}))^2}{\gamma_1^4(\gamma_1^2-1)} + \frac{2 L(\gamma_1^{-1}) L''(\gamma_1^{-1})}{\gamma_1^4 (\gamma_1^2 - 1)} - \frac{2L^2(\gamma_1^{-1})}{(\gamma_1^2 - 1)^2}+\frac{4(2\gamma_1^2-1)L(\gamma_1^{-1})L'(\gamma_1^{-1})}{\gamma_1^3(\gamma_1^2-1)^2}\\
      &~~~~ + \frac{4L(\gamma_1^{-1})L'(\gamma_1^{-1})}{\gamma_1(\gamma_1^2 - 1)^2}-\frac{2L(\gamma^{-1})L''(\gamma_1^{-1})}{\gamma_1^4(\gamma_1\gamma-1)}- \frac{2(3\gamma_1\gamma-2)L(\gamma^{-1})L'(\gamma_1^{-1})}{\gamma_1^3(\gamma_1\gamma-1)^2}\\
      &~~~~-\frac{2\gamma L(\gamma^{-1})L'(\gamma_1^{-1})}{\gamma_1^2(\gamma_1\gamma-1)^2} - \frac{4\gamma^2 L(\gamma^{-1})L(\gamma_1^{-1})}{(\gamma_1\gamma-1)^3}+\frac{8 \gamma_1^2 L^2(\gamma_1^{-1})}{(\gamma_1^2 - 1)^3},
   \end{aligned}
\end{equation*}
which implies that
\begin{equation}\label{g-2}
   \frac{\partial^2g}{\partial \gamma_1^2}(\gamma,\gamma)=\frac{2\left(L^\prime(\gamma^{-1})\right)^2}{\gamma^4(\gamma^2-1)} +\frac{4L(\gamma^{-1}) L'(\gamma^{-1})}{\gamma(\gamma^2-1)^2} +\frac{2(\gamma^2 + 1)\left(L(\gamma^{-1})\right)^2}{(\gamma^2 - 1)^3}.
\end{equation}
Therefore, combining \eqref{f}, \eqref{taylor-g}, \eqref{g-0}, \eqref{g-1} and \eqref{g-2} yields
\begin{equation*}
\begin{aligned}
   \lim_{\gamma_1 \to \gamma} f(\gamma_1, \gamma) &= \frac{(\gamma^2 - 1)^4}{2 L^4 (\gamma^{-1})} \frac{\partial^2 g}{\partial \gamma_1^2} (\gamma, \gamma)\\
   &=\frac{(\gamma^2 - 1)^3}{\gamma^4} 
\frac{\left(L'(\gamma^{-1})\right)^2}{L^4(\gamma^{-1})} + \frac{2 (\gamma^2 - 1)^2}{\gamma} \frac{L'(\gamma^{-1})}{L^3(\gamma^{-1})} + \frac{\gamma^4 - 1}{L^2(\gamma^{-1})},
\end{aligned}
\end{equation*}
from which (\ref{power-eqn-equivalence-case-c}) follows immediately.






\newpage

\bibliographystyle{ieeetr}
\bibliography{ArP_SK2_v6}

\end{document}
